{\color{black}
\section{Chapter 1. Introduction to causality}
\sourcepages{Chapter1-23.pdf, p. 1}
$T_i$ is treatment intake; $Y_i$ is the observed outcome. The notebook contrasts ``reality (group assignment)'' with ``truth.''
\[
\mathrm{ATE}=E[Y_1-Y_0],\qquad \mathrm{ATT}=E[Y_1-Y_0\mid T=1].
\]
ATE is the average treatment effect; ATT concerns those who receive treatment. The causal quantity is the potential-outcome contrast. The observed association decomposes as
\[
\underbrace{E[Y\mid T=1]-E[Y\mid T=0]}_{\inkblue{\text{association}}}
=\underbrace{E[Y_1-Y_0\mid T=1]}_{\inkblue{\mathrm{ATT}}}
+\underbrace{E[Y_0\mid T=1]-E[Y_0\mid T=0]}_{\inkblue{\text{bias}}}.
\]
\begin{inkblock}{blue}
\textbf{My annotation.} When bias is zero, association represents ATT; the arrows then connect ATT to $E[Y_1-Y_0]$, labeled ``causality.''
\end{inkblock}

\par\Needspace{9\baselineskip}\smallskip\textbf{Original pointing relations.}\par\smallskip
\notearrow{blue}{Bias $=0$: association represents ATT}{$\mathrm{ATT}$}
\notearrow{blue}{Can be converted to causality (source arrow)}{$E[Y_1-Y_0]$; see the qualification below}
\begin{reviewnote}{qualification}{B1-004}{Zero selection bias identifies ATT, not automatically ATE}{Chapter1-23.pdf, pp. 1}
The arrow is present in the original blue ink. Zero $E[Y_0\mid T=1]-E[Y_0\mid T=0]$ gives association $=\mathrm{ATT}$. To replace ATT by ATE also require $E[Y_1-Y_0\mid T=1]=E[Y_1-Y_0]$; random assignment of both potential outcomes is sufficient. Indeed, $\mathrm{ATE}=P(T=1)\mathrm{ATT}+P(T=0)\mathrm{ATU}$.
\par\smallskip\textit{Reference:} \href{https://matheusfacure.github.io/python-causality-handbook/01-Introduction-To-Causality.html}{Handbook, Chapter 1}.
\end{reviewnote}
\section{Chapter 2. Randomized experiments}
\sourcepages{Chapter1-23.pdf, p. 1}
Randomized experiments provide a way to remove bias. The sketch separates the link from existing circumstances to potential outcomes (the link to eliminate) from the treatment-to-observed-outcome link (the link to study). An RCT supplies
\begin{inkblock}{blue}
\[
(Y_0,Y_1)\perp T,\qquad E[Y_0\mid T=0]=E[Y_0\mid T=1]=E[Y_0],
\]
so the selection-bias term is zero.
\end{inkblock}

\par\Needspace{9\baselineskip}\smallskip\textbf{Original pointing relations.}\par\smallskip
\notearrow{blue}{RCT supplies the condition}{$(Y_0,Y_1)\perp T$}
\notearrow{blue}{Need to eliminate this link}{Existing circumstances $\longrightarrow$ potential outcomes}
\notearrow{blue}{Need to study this link}{Treatment $\longrightarrow$ observed outcome}
\section{Chapter 3. A review of statistics}
\sourcepages{Chapter1-23.pdf, p. 2}
The standard error of an estimate is written as $\mathrm{SE}=\hat\sigma/\sqrt n$, with
\[\hat\sigma^2=\frac{1}{N-1}\sum_i(x_i-\bar x)^2.\]
Smaller sample size $\Rightarrow$ larger standard error $\Rightarrow$ wider confidence interval. The null hypothesis is that the treatment and control distributions have the same mean. Even if their individual confidence intervals overlap, their mean difference may still be statistically different from zero. \inkblue{The blue annotations distinguish the 95\% intervals for the two group means from the difference between those means.}

The notebook writes the $p$-value as $P(\text{data}\mid H_0)$: under $H_0$, the probability of the observed data or more extreme data (with a Type I error reminder). Confidence level is described as the probability that the estimate and the population parameter differ by no more than a specified margin. Higher confidence $\Rightarrow$ a wider interval $\Rightarrow$ lower precision. The significance level is the probability of a Type I error.

\textbf{My intuition.} If the $p$-value is very small, then the observed data would be very unlikely under $H_0$, yet they occurred, suggesting that $H_0$ is false. \inkblue{Compare the sample's $p$-value with the significance threshold chosen theoretically.}

\par\Needspace{9\baselineskip}\smallskip\textbf{Original pointing relations.}\par\smallskip
\notearrow{black}{$\mathrm{SE}=\hat\sigma/\sqrt n$}{$\hat\sigma^2=(N-1)^{-1}\sum_i(x_i-\bar x)^2$}
\notearrow{blue}{Each confidence interval}{95\% interval for the mean in the $T=1$ or $T=0$ group}
\notearrow{blue}{The difference can still be significant}{Difference between the two group means}
\begin{reviewnote}{qualification}{B1-005}{Sampling probability and confidence coverage}{Chapter1-23.pdf, pp. 2}
The scan's shorthand needs its sampling interpretation. For an i.i.d. sample mean, $\widehat{\mathrm{SE}}(\bar X)=s/\sqrt n$; this is not every estimator's SE. A $p$-value is $P_{H_0}\{S\text{ at least as extreme as }S_{\rm obs}\}$, not $P(H_0\mid\text{data})$. A 95\% confidence procedure has repeated-sampling coverage $P_\theta\{\theta\in C(X)\}=0.95$; a small $p$-value is evidence against the null, not a logical proof of its falsity.
\par\smallskip\textit{Reference:} \href{https://matheusfacure.github.io/python-causality-handbook/03-Stats-Review-The-Most-Dangerous-Equation.html}{Handbook, Chapter 3}.
\end{reviewnote}
\section{Chapter 4. Graphical causal models}
\sourcepages{Chapter1-23.pdf, p. 3}
$A\to B$ means that $A$ affects $B$; dependence propagates along the arrows.
\begin{enumerate}
\item Chain: $A\to B\to C$. Marginally $A\not\perp C$; conditioning on the known intermediate variable $B$ blocks the path: $A\perp C\mid B$.
\item Fork: $A\leftarrow C\to B$. Without conditioning, $A\not\perp B$; conditioning on $C$ gives $A\perp B\mid C$.
\item Collider: $A\to C\leftarrow B$. Marginally $A\perp B$; conditioning on $C$ gives $A\not\perp B\mid C$.
\end{enumerate}
Confounding bias may involve proxy confounders and variables that cannot be directly measured.

\textbf{My summary.} Confounding bias arises from an uncontrolled open path, illustrated by $X\leftarrow W\to Y$ together with $X\to Y$. Selection bias comes from extra conditioning: conditioning on a collider opens a path, or conditioning on an intermediate variable blocks part of the causal path. The example labels $X$ as education, $Y$ as wages, and $W$ as investment. The collider sketch has $X\to W_2\leftarrow Y$, $W_2\to W_1$, and $X\to Y$, with ``investment'' written alongside $W_1$; the mediator sketch has education feeding $W$, then wages, alongside a direct education-to-wages arrow. The note says controlling $W$ also has a negative effect on this analysis.

\begin{center}
\begin{tikzpicture}[>=Stealth, font=\small, line width=.65pt]

\node (a) at (0,0) {$A$}; \node (b) at (1.5,0) {$B$}; \node (c) at (3,0) {$C$}; \draw[->] (a)--(b); \draw[->] (b)--(c);
\node (d) at (5,0) {$A$}; \node (e) at (6,1) {$C$}; \node (f) at (7,0) {$B$}; \draw[->] (e)--(d); \draw[->] (e)--(f);
\node (g) at (9,1) {$A$}; \node (h) at (11,1) {$B$}; \node (i) at (10,0) {$C$}; \draw[->] (g)--(i); \draw[->] (h)--(i);
\node (w) at (0,-2.1) {$W$}; \node (x) at (1.5,-1.4) {$X$}; \node (y) at (1.5,-2.8) {$Y$}; \draw[->] (w)--(x);\draw[->] (w)--(y);\draw[->] (x)--(y);
\node (x2) at (5.5,-1.4) {$X$};\node (y2) at (5.5,-2.8) {$Y$};\node (w2) at (4.2,-2.1) {$W_2$};\node (w1) at (2.9,-2.1) {$W_1$};\draw[->] (x2)--(y2);\draw[->] (x2)--(w2);\draw[->] (y2)--(w2);\draw[->] (w2)--(w1);
\node (x3) at (10,-1.4) {$X$};\node (y3) at (10,-2.8) {$Y$};\node (w3) at (8.7,-2.1) {$W$};\draw[->] (x3)--(y3);\draw[->] (x3)--(w3);\draw[->] (w3)--(y3);

\end{tikzpicture}
\par{\small Source p. 3: chain, fork, collider; followed by confounding and the two bad-control graphs.}
\end{center}
\par\Needspace{9\baselineskip}\smallskip\textbf{Original pointing relations.}\par\smallskip
\notearrow{black}{Extra conditioning on a common effect opens a path}{Collider $W_2$, or its descendant $W_1$}
\notearrow{black}{Extra conditioning blocks part of the causal path}{Mediator $W$ in $X\to W\to Y$}
\begin{reviewnote}{qualification}{B1-006}{Graphical separation and dependence}{Chapter1-23.pdf, pp. 3}
These are the source's three-node path rules. Under the causal Markov property, a blocked path implies the stated independence when no other path is open. An open path permits dependence; it does not force it without faithfulness (effects can cancel). Adjusting for a mediator blocks part of a total effect; that is overcontrol, whereas conditioning on a collider can create selection bias.
\par\smallskip\textit{Reference:} \href{https://matheusfacure.github.io/python-causality-handbook/04-Graphical-Causal-Models.html}{Handbook, Chapter 4}.
\end{reviewnote}
\section{Chapter 5. The usefulness of linear regression}
\sourcepages{Chapter1-23.pdf, p. 4}
For individual $i$, write $Y_{1i}=Y_{0i}+k$ and estimate $k$. In a single regression,
\[Y_i=\beta_0+kT_i+u_i,\qquad \hat\beta_1=\frac{\operatorname{Cov}(Y_i,T_i)}{\operatorname{Var}(T_i)}.\]
The drawn line joins the intercept at $T=0$ to intercept plus $k$ at $T=1$ (a negative slope in the sketch). When $T$ is randomly assigned, its regression coefficient is ATE.

Multiple regression addresses the many factors that affect $T_i$:
\[y_i=\beta_0+kT_i+\beta_1X_{1i}+\cdots+\beta_KX_{Ki}+u_i.\]
\inkblue{The $X$ terms are auxiliary covariates.} Then
\[k=\frac{\operatorname{Cov}(Y_i,\widetilde T_i)}{\operatorname{Var}(\widetilde T_i)}.\]
\begin{inkblock}{blue}
\textbf{My annotation.} First remove the part explained by the $X$ variables and take the residual $\widetilde T$: the remaining part cannot be explained by $X$.
\end{inkblock}

If a multiple-regression model is needed but a simple regression omits a variable, consider the correct model
\[\mathrm{Wage}_i=\alpha+k\,\mathrm{Educ}_i+A_i'\beta+u_i,\]
where $A_i$ is ability. The omitted-variable calculation is
\[\frac{\operatorname{Cov}(\mathrm{Wage}_i,\mathrm{Educ}_i)}{\operatorname{Var}(\mathrm{Educ}_i)}=k+\beta'\delta_{\mathrm{Ability}}.\]
\inkblue{$\delta_A$ is the coefficient vector from regressing ability on education.} \inkred{The extra term is omitted variable bias (OVB).} If all confounders are included, there is no OVB. \textbf{Margin:} an RCT breaks the relationship; regression controls it and holds it fixed.

\begin{center}
\begin{tikzpicture}[>=Stealth, font=\small, line width=.65pt]

\draw[->] (0,0)--(4,0) node[right] {$T$};\draw[->] (0,0)--(0,2.5) node[above] {$Y$};
\draw (0,2)--(3,1);\draw[dashed] (0,1) node[left] {$\beta_0+k$}--(3,1)--(3,0) node[below] {$1$};\node[left] at (0,2) {$\beta_0$};\node[below] at (0,0) {$0$};

\end{tikzpicture}
\par{\small Source p. 4: binary-treatment regression (negative slope as drawn).}
\end{center}
\par\Needspace{9\baselineskip}\smallskip\textbf{Original pointing relations.}\par\smallskip
\notearrow{blue}{Use $X$ to predict first, then take the unexplained part}{$\widetilde T=T-\widehat{E[T\mid X]}$ (linear projection here)}
\notearrow{blue}{The omitted term is ability}{$A_i^{\prime}\beta$}
\notearrow{blue}{Ability regressed on education}{$\delta_A$}
\notearrow{red}{OVB: omitted variable bias}{$\beta^{\prime}\delta_A$}
\begin{reviewnote}{original}{B1-007}{Direction of the residual regression}{Chapter1-23.pdf, pp. 4}
This reversal is in the original black sentence, and the blue annotation already points toward the repair. Let $P_X$ include the intercept and set $\widetilde T=(I-P_X)T$. Then $\hat k=\widetilde T'Y/(\widetilde T'\widetilde T)$. Regressing $X$ on $T$ instead does not produce this residual.
\par\smallskip\textit{Reference:} \href{https://matheusfacure.github.io/python-causality-handbook/05-The-Unreasonable-Effectiveness-of-Linear-Regression.html}{Handbook, Chapter 5}.
\end{reviewnote}

\begin{reviewnote}{qualification}{B1-008}{Regression adjustment still needs identification}{Chapter1-23.pdf, pp. 4}
Including all confounders is sufficient only with consistency, overlap, conditional exchangeability and an adequate adjustment specification. The displayed OVB identity uses $\operatorname{Cov}(\mathrm{Educ},u)=0$. An additive linear regression can remain misspecified even with all relevant covariates measured; adding them does not automatically identify the unweighted ATE under heterogeneous effects.
\par\smallskip\textit{Reference:} \href{https://matheusfacure.github.io/python-causality-handbook/05-The-Unreasonable-Effectiveness-of-Linear-Regression.html}{Handbook, Chapter 5}.
\end{reviewnote}
\section{Chapter 6. Grouping and dummy regression}
\sourcepages{Chapter1-23.pdf, p. 5}
Heteroskedasticity means that outcome variance is not constant across the values of the covariates. With finer groups, the number of observations falls; grouping reduces variation, but coefficient standard errors can rise and $t$ statistics become smaller. Group-specific regressions should be weighted by within-group treatment variance.

For a dummy treatment (the note gives ``whether 12 years'' as the example),
\[\mathrm{wage}_i=\beta_0+\beta_1T_i+e_i.\]
At $T=0$ the expectation is the intercept; at $T=1$ it is intercept plus $\beta_1$. With continuous IQ,
\[\mathrm{wage}_i=\beta_0+\beta_1T_i+\beta_2\mathrm{IQ}_i+\beta_3\mathrm{IQ}_iT_i+e_i.\]
\begin{inkblock}{blue}
\begin{itemize}
\item $\beta_0$: the expectation at $T=0$, IQ $=0$.
\item $\beta_1$: the effect of one unit of $T$ when IQ is zero (all other terms are zero).
\item $\beta_2$: the effect of one unit of IQ when $T=0$.
\item $\beta_3$: IQ also has a return when $T=1$, making that group's regression slope larger when $\beta_3>0$.
\end{itemize}
\end{inkblock}
To analyze the coefficient on $T_i$, calculate each group's contribution to the overall coefficient: group by years of education (12-year groups in the note) and compute the corresponding categorical education/treatment contrast $C(\mathrm{educ})[T_i]$.

With more covariates, the target coefficient is a weighted average of the effect in each dummy-variable subgroup:
\[E\!\left[\{E[Y_i\mid T=1,G_i]-E[Y_i\mid T=0,G_i]\}\,w(G_i)\right],\qquad
w(G_i)\propto\operatorname{Var}(T_i\mid G_i).
\]
\inkblue{The weights depend on the distribution of $T_i$ within a subgroup. The margin says to find both the education and IQ-bin treatment contrasts.}

\par\Needspace{9\baselineskip}\smallskip\textbf{Original pointing relations.}\par\smallskip
\notearrow{black}{The binary treatment changes from 0 to 1}{$E[Y\mid T=1]-E[Y\mid T=0]=\beta_1$ in this regression}
\notearrow{blue}{IQ also changes the slope for $T=1$}{Interaction coefficient $\beta_3$}
\begin{reviewnote}{qualification}{B1-009}{Normalize the subgroup regression weights}{Chapter1-23.pdf, pp. 5}
With a separate intercept for every group and a common treatment coefficient, $\beta=\sum_g P(G=g)v_g\tau_g/\sum_gP(G=g)v_g$, where $v_g=\operatorname{Var}(T\mid G=g)$. Equivalently, the scan's outer expectation uses $w(g)=v_g/E[v_G]$. This is generally a variance-weighted effect, not the ordinary ATE; causal interpretation additionally needs valid within-group adjustment.
\par\smallskip\textit{Reference:} \href{https://matheusfacure.github.io/python-causality-handbook/06-Grouped-and-Dummy-Regression.html}{Handbook, Chapter 6}.
\end{reviewnote}
\section{Chapter 7. Beyond controlling for confounders}
\sourcepages{Chapter1-23.pdf, p. 6}
This chapter concerns the effects of controlling for given variables. In $X\to Y\leftarrow T$, controlling $T$ when studying $X\to Y$ can make the effect of $X$ clearer and reduce the variance of its estimate. A second sketch has $B\to A\to X\to Y$ and $B\to Y$: $A$ is related only to $X$, while $B$ is related to both $X$ and $Y$. \inkblue{The annotation says that adding $A$, related only to $X$, increases the variance (the handwritten outcome label is $Y$).}
\[\hat\sigma^2=\frac{1}{n-2}\sum_i(y_i-\hat y_i)^2,\qquad
\operatorname{Var}(\hat\beta_2)=\frac{\hat\sigma^2}{\sum_i(x_i-\bar x)^2}.\]
The adjacent black prose says ``standard error is inversely proportional to the variance of $X$.'' \inkblue{Larger sample variation gives a smaller standard error.}

Confounding is the bias from failing to control for a common cause. Selection bias occurs when controlling for a common effect or a variable on the path from cause to effect.

\paragraph{Bad controls and participation.}
For a treatment that affects whether the outcome is zero, an RCT makes the mean difference causal. The note decomposes it into a participation effect and an effect among positive outcomes:
\begin{align*}
E[Y\mid T=1]-E[Y\mid T=0]
={}&\{P(Y>0\mid T=1)-P(Y>0\mid T=0)\}E[Y\mid Y>0,T=1]\\
&+\{E[Y\mid Y>0,T=1]-E[Y\mid Y>0,T=0]\}P(Y>0\mid T=0).
\end{align*}
\inkblue{The margin identifies the first part as the effect on the probability of a positive outcome, and distinguishes customers' participation from spending ability.} Restricting analysis to participants gives
\begin{inkblock}{blue}
\begin{align*}
E[Y_{i1}\mid Y_{i1}>0]-E[Y_{i0}\mid Y_{i0}>0]
={}&\underbrace{E[Y_{i1}-Y_{i0}\mid Y_{i1}>0]}_{\text{causal effect among participants}}\\
&+\underbrace{E[Y_{i0}\mid Y_{i1}>0]-E[Y_{i0}\mid Y_{i0}>0]}_{\text{selection bias}}.
\end{align*}
The last term compares participants and nonparticipants [the end of the annotation is clipped]. \end{inkblock}
The accompanying graph has $X\to T_1\to Y\to T_2$, and $X\to T_2$. \inkred{$T_1,T_2$ are marked in red as bad selections/controls: ``do not control!''}

\begin{center}
\begin{tikzpicture}[>=Stealth, font=\small, line width=.65pt]

\node[draw,circle] (x) at (0,1) {$X$};\node (t) at (2,1) {$T$};\node[draw,circle] (y) at (1,0) {$Y$};\draw[->] (x)--(y);\draw[->] (t)--(y);
\node (b) at (5.5,1.5) {$B$};\node (a) at (4.2,.8) {$A$};\node[draw,circle] (xx) at (4.2,-.2) {$X$};\node[draw,circle] (yy) at (5.5,-1) {$Y$};\draw[->] (b)--(a);\draw[->] (a)--(xx);\draw[->] (xx)--(yy);\draw[->] (b)--(yy);
\node[draw,circle] (xxx) at (9,1.5) {$X$};\node[draw=red,circle] (t1) at (8,.5) {$T_1$};\node[draw,circle] (yyy) at (8.5,-.6) {$Y$};\node[draw=red,circle] (t2) at (10,-1) {$T_2$};\draw[->] (xxx)--(t1);\draw[->] (t1)--(yyy);\draw[->] (yyy)--(t2);\draw[->] (xxx)--(t2);

\end{tikzpicture}
\par{\small Source p. 6: good precision control, treatment predictor, and bad controls (red circles).}
\end{center}
\par\Needspace{9\baselineskip}\smallskip\textbf{Original pointing relations.}\par\smallskip
\notearrow{blue}{Only related to $X$; adding it increases variance (source wording)}{$A$ in the second graph}
\notearrow{blue}{Related to both $X$ and $Y$}{$B$ in the second graph}
\notearrow{red}{Bad controls: do not condition on these}{$T_1$ (mediator) and $T_2$ (collider) in the lower graph}
\notearrow{blue}{Participation effect}{$P(Y>0\mid T=1)-P(Y>0\mid T=0)$}
\notearrow{blue}{Selection bias}{$E[Y_0\mid Y_1>0]-E[Y_0\mid Y_0>0]$}
\begin{reviewnote}{qualification}{B1-010}{Precision, bad controls, and participation}{Chapter1-23.pdf, pp. 6}
For homoskedastic simple OLS, $\mathrm{SE}(\hat\beta)=\sigma/\sqrt{\sum_i(x_i-\bar x)^2}$: SE is inversely proportional to the square root of treatment variation. In multiple regression use residual treatment and the appropriate residual degrees of freedom. The variance intuition concerns the treatment-effect estimator, not literally a change in $\operatorname{Var}(Y)$. The participation decomposition assumes $Y\ge0$ (or that zero and positive outcomes exhaust the support); its two algebraic components are not separately causal effects without further assumptions.
\par\smallskip\textit{Reference:} \href{https://matheusfacure.github.io/python-causality-handbook/07-Beyond-Confounders.html}{Handbook, Chapter 7}.
\end{reviewnote}
\section{Chapter 8. Instrumental variables}
\sourcepages{Chapter1-23.pdf, pp. 7--8}
Suppose
\[Y_i=\beta_0+kT_i+\beta W_i+u_i,\]
where $T$ is treatment and $W$ is a confounder. When $W_i$ is unavailable, fit $Y_i=\beta_0+kT_i+v_i$, with $v_i=\beta W_i+u_i$. Then $\operatorname{Cov}(T,v)\ne0$, so the estimate of $k$ is biased. \inkblue{The unobserved term carries information linking $T$ and $Y$.}

Introduce an instrument $Z$ related to $T$ and affecting $Y$ only through $T$, with $\operatorname{Cov}(Z,v)=0$. The graph is $Z\to T\to Y$, $W\to T$, $W\to Y$. Thus
\begin{align*}
\operatorname{Cov}(Z,Y)&=\operatorname{Cov}(Z,\beta_0+kT+v)\\
&=k\operatorname{Cov}(Z,T)+\operatorname{Cov}(Z,v)=k\operatorname{Cov}(Z,T),\\
k&=\frac{\operatorname{Cov}(Z,Y)}{\operatorname{Cov}(Z,T)}
=\frac{\operatorname{Cov}(Y,Z)/\operatorname{Var}(Z)}{\operatorname{Cov}(T,Z)/\operatorname{Var}(Z)}.
\end{align*}
Both numerator and denominator are regression coefficients. \inkblue{The numerator is the reduced-form coefficient: the effect of $Z$ on $Y$, operating only through $T$.}

\textbf{My calculus intuition.}
\[k=\frac{dY/dZ}{dT/dZ}=\frac{dY}{dZ}\frac{dZ}{dT}=\frac{dY}{dT}.\]
For a binary instrument, the Wald estimator is
\[k=\frac{E[Y\mid Z=1]-E[Y\mid Z=0]}{E[T\mid Z=1]-E[T\mid Z=0]}.\]
The note describes this as turning the change from $Z\to Y$ back into the effect $T\to Y$.

\paragraph{Using an instrument.}
The example labels $Z_0$ ``qob'' (quarter of birth), $T$ education, $W$ ability, and $Y$ wages; additional variables are $Z_1,Z_2$. The proposed requirements are (1) $\operatorname{Cov}(Z,T)\ne0$, which can be checked directly in the data, and (2) $Y\perp Z\mid T$, followed by question marks and ``$Z$ affects $Y$ only through $T$.'' \inkred{\textbf{Author's red correction:} ``Not accurate; corrected in Chapter 09.''} This conditional-independence statement is therefore retained as a superseded note, not an accepted assumption.

Regress $T\sim Z_0$ to check the first-stage relationship, then $Y\sim Z_0$ to examine the reduced form. Method 1 takes the ratio of these regression coefficients. Method 2 is 2SLS:
\begin{enumerate}
\item Regress $T$ on $Z_0,Z_1,Z_2$.
\item Regress $Y$ on fitted $T$ and $Z_1,Z_2$; the sketch substitutes $(Z_0,Z_1,Z_2)$ for fitted $T$.
\end{enumerate}
\inkblue{The coding note says that one can specify $T\sim Z_0$ because the other terms are automatically added in stage 2.} The OLS comparison is written $Y\sim T+C(Z_0,Z_1,Z_2)$.

With a weak correlation between $T$ and $Z$, estimates differ between OLS and 2SLS, and standard errors are large [the last parenthetical word is unclear]. 2SLS has finite-sample bias toward OLS, with the same direction as OLS bias; 2SLS is consistent while OLS is not. Bias increases with the number of instruments: with too many IVs, 2SLS increasingly resembles OLS.

\begin{center}
\begin{tikzpicture}[>=Stealth, font=\small, line width=.65pt]

\node (z) at (0,1) {$Z_0$};\node[blue] at (0,1.55) {quarter of birth};\node[draw,circle] (t) at (3,1) {$T$};\node[blue] at (3,1.55) {education};\node (w) at (6,1) {$W$};\node[blue] at (6,1.55) {ability};\node[draw,circle] (y) at (3,-1) {$Y$};\node[blue] at (3,-1.6) {wages};\node (z1) at (0,0) {$Z_1$};\node (z2) at (0,-1) {$Z_2$};\draw[->] (z)--(t);\draw[->] (z1)--(t);\draw[->] (z2)--(t);\draw[->] (w)--(t);\draw[->] (w)--(y);\draw[->] (t)--(y);

\end{tikzpicture}
\par{\small Source pp. 7--8: instrument, treatment, confounder and outcome.}
\end{center}
\par\Needspace{9\baselineskip}\smallskip\textbf{Original pointing relations.}\par\smallskip
\notearrow{blue}{Information connecting $T$ and $Y$}{$v_i=\beta W_i+u_i$}
\notearrow{blue}{Reduced form: $Z$ affects $Y$ through $T$}{$\operatorname{Cov}(Y,Z)/\operatorname{Var}(Z)$}
\notearrow{red}{Not accurate; corrected in Chapter 09}{The proposed condition $Y\perp Z\mid T$}
\begin{reviewnote}{original}{B1-011}{Exclusion is not conditional observed-outcome independence}{Chapter1-23.pdf, pp. 8}
The author's red correction is appropriate. Exclusion is $Y(t,z)=Y(t)$, together with instrument exogeneity; it is not generally $Y\perp Z\mid T$. In the drawn graph, conditioning on collider $T$ opens $Z\to T\leftarrow W\to Y$. Hence a valid instrument can remain associated with observed $Y$ after conditioning on treatment.
\par\smallskip\textit{Reference:} \href{https://matheusfacure.github.io/python-causality-handbook/09-Non-Compliance-and-LATE.html}{Handbook, Chapter 9}.
\end{reviewnote}

\begin{reviewnote}{qualification}{B1-012}{IV ratios and finite-sample claims}{Chapter1-23.pdf, pp. 7, 8}
The covariance ratio follows from $\operatorname{Cov}(Z,v)=0$ and a nonzero first stage in the stated constant-slope model. The derivative cancellation is an intuition for a differentiable single pathway, not a general proof for stochastic heterogeneous IV models. Ordinary 2SLS consistency requires valid sufficiently strong instruments and a fixed-dimensional regular asymptotic regime; weak or many instruments invalidate an unconditional consistency claim. Treat $Z_1,Z_2$ here as included exogenous controls, present in both stages.
\par\smallskip\textit{Reference:} \href{https://matheusfacure.github.io/python-causality-handbook/08-Instrumental-Variables.html}{Handbook, Chapter 8}.
\end{reviewnote}
\section{Chapter 9. Noncompliance and local average treatment effects}
\sourcepages{Chapter1-23.pdf, pp. 9--10}
Compliers follow assignment; never-takers never receive treatment; always-takers always receive it; defiers do the opposite. In $Z\to T\to Y$ with $W\to T$ and $W\to Y$, never-takers and always-takers weaken the first stage. \inkblue{The blue intuition says that, for example, never-takers pull the reduced-form outcome difference toward the control condition.}

Potential outcomes now have two arguments: $Y_i(t,z)$; the written example is $Y_i(1,0)$ for $T_i=1,Z_i=0$. Potential treatment $T_{0i}$ corresponds to $Z_i=0$.
\paragraph{IV assumptions, as recorded.}
\begin{enumerate}
\item $(T_{0i},T_{1i})\perp Z_i$ and $(Y_i(T_{1i},1),Y_i(T_{0i},0))\perp Z_i$. $Z$ is unrelated to potential treatment; people in different instrument groups are comparable.
\item Exclusion: $Y_i(1,0)=Y_i(1,1)=Y_{i1}$ and $Y_i(0,0)=Y_i(0,1)=Y_{i0}$. The instrument affects the outcome only via treatment.
\item $E[T_{1i}-T_{0i}]\ne0$: the first stage exists; the instrument really affects treatment.
\item The scan writes $T_{1i}>T_{0i}$ and says switching on the instrument increases the probability of treatment.
\end{enumerate}

\paragraph{Deriving the Wald estimand.}
The scan initially labels the ratio ``ATE'':
\[\frac{E[Y\mid Z=1]-E[Y\mid Z=0]}{E[T\mid Z=1]-E[T\mid Z=0]}.\]
Using exclusion and random assignment,
\begin{align*}
E[Y_i\mid Z_i=1]&=E[Y_{i0}+T_{i1}(Y_{i1}-Y_{i0})\mid Z_i=1]\\
&=E[Y_{i0}+T_{i1}(Y_{i1}-Y_{i0})],\\
E[Y_i\mid Z_i=0]&=E[Y_{i0}+T_{i0}(Y_{i1}-Y_{i0})].
\end{align*}
The numerator is
\begin{align*}
E[(Y_{i1}-Y_{i0})(T_{i1}-T_{i0})]
=E[Y_{i1}-Y_{i0}\mid T_{i1}>T_{i0}]P(T_{i1}>T_{i0}),
\end{align*}
and the denominator is $E[T_{i1}-T_{i0}]=P(T_{i1}>T_{i0})$. Their ratio is therefore
\[E[Y_{i1}-Y_{i0}\mid T_{i1}>T_{i0}].\]
\inkred{\textbf{Author's red annotation:} this is the condition attached to that average treatment effect: compliance.} IV identifies the treatment effect only for compliers. Compliers have $T_{i1}>T_{i0}$; never-takers have $T_{i1}=T_{i0}=0$; always-takers have $T_{i1}=T_{i0}=1$. Defiers run counter to assignment.

\par\Needspace{9\baselineskip}\smallskip\textbf{Original pointing relations.}\par\smallskip
\notearrow{blue}{Never-/always-takers weaken the first stage}{$Z\longrightarrow T$}
\notearrow{black}{Instrument groups are comparable}{Instrument exogeneity assumptions}
\notearrow{black}{Instrument affects outcome only through treatment}{$Y_i(t,0)=Y_i(t,1)$}
\notearrow{red}{This condition means compliance}{$T_{i1}>T_{i0}$ in the final conditional expectation}
\begin{reviewnote}{original}{B1-013}{Non-strict monotonicity and the LATE label}{Chapter1-23.pdf, pp. 9, 10, 23}
The original assumptions use $T_1>T_0$, but the subsequent types include $T_1=T_0$. Replace the population condition by $T_1\ge T_0$ almost surely and $P(T_1>T_0)>0$. Under joint instrument exogeneity and exclusion, the numerator is $P(C)E[Y_1-Y_0\mid C]$ and the denominator $P(C)$, with $C=\{T_1>T_0\}$. The ratio is LATE, not population ATE. The same non-strict condition applies to the fuzzy-RD discussion on p.~23.
\par\smallskip\textit{Reference:} \href{https://matheusfacure.github.io/python-causality-handbook/09-Non-Compliance-and-LATE.html}{Handbook, Chapter 9}; \href{https://www.nber.org/papers/t0118}{Imbens--Angrist, LATE identification}.
\end{reviewnote}
\section{Chapter 10. Matching}
\sourcepages{Chapter1-23.pdf, pp. 11--12}
When grouping on confounders to estimate ATE, weight the treatment effect in each group. Weighting by sample size pulls the estimate toward larger groups, even though their treatment variation may be small; this is marked ``nonparametric estimation.'' \inkblue{Weighting by variance gives higher-variance groups greater weight; this is linked to linear, parametric regression and its emphasis on high-variance covariates.}

\paragraph{Subclassification.}
If $X$ creates confounding, group by equal $X$ values. Under $(Y_0,Y_1)\perp T\mid X$,
\begin{align*}
\mathrm{ATE}&=\int\{E[Y\mid X=x,T=1]-E[Y\mid X=x,T=0]\}\,dP(x),\\
\widehat{\mathrm{ATE}}&=\sum_{k=1}^{K}(\bar Y_{k1}-\bar Y_{k0})\frac{N_k}{N}.
\end{align*}
\inkblue{The blue annotation describes averaging over individual covariate values $\{x_1,\ldots,x_K\}$.}

\paragraph{Matching estimator.}
Use distances between observations to define similarity, then compare similar or nearby individuals receiving different treatments. \inkred{K-nearest neighbors is noted.} \inkblue{Normalize a coordinate when defining distance: $(x_i-\bar x)/\sqrt{\operatorname{Var}(x)}$.}
\[\widehat{\mathrm{ATE}}=\frac1N\sum_i(2T_i-1)(Y_i-Y_{j(i)}).\]
\inkblue{$j(i)$ is the most similar observation from the other treatment group.} \inkred{The factor $2T_i-1$ changes the sign according to which treatment group $i$ belongs to.}

\paragraph{Matching bias and correction.}
For the treated group, consider
\[\widehat{\mathrm{ATET}}=\frac1{N_1}\sum_{i:T_i=1}(Y_i-Y_{j(i)}),\qquad
\inkblue{\mu_0(x)=E[Y\mid X=x,T=0]}.\]
The centering discrepancy is written
\[E[\sqrt{N_1}(\widehat{\mathrm{ATET}}-\mathrm{ATET})]
=E[\sqrt{N_1}\{\mu_0(X_i)-\mu_0(X_{j(i)})\}],\]
and need not converge to zero. \inkblue{Multiplying the discrepancy by $\sqrt{N_1}$ magnifies it.} \inkblue{The blue explanation contrasts $X_i$'s untreated potential outcome (counterfactual because $i$ is treated) with $X_j$'s observed untreated outcome; $X_i\approx X_j$ is meant to make their untreated outcomes close.} Correct the estimator so that it converges:
\[\widehat{\mathrm{ATET}}_{\mathrm{bc}}=
\frac1{N_1}\sum_{i:T_i=1}\left[(Y_i-Y_{j(i)})-
\{\hat\mu_0(X_i)-\hat\mu_0(X_{j(i)})\}\right].\]
\inkblue{$\hat\mu_0$ is estimated, for example by least squares.} A local LS correction does not alter the matching estimator's original aim.

\textbf{The curse of dimensionality.} More covariates mean larger matching distances: each cell contains fewer observations.

\par\Needspace{9\baselineskip}\smallskip\textbf{Original pointing relations.}\par\smallskip
\notearrow{blue}{Average over the covariate values}{$\int \tau(x)\,dP(x)$}
\notearrow{red}{Sign coding by treatment group}{$2T_i-1$}
\notearrow{blue}{Most similar member of the other treatment group}{$j(i)$}
\notearrow{blue}{Counterfactual untreated mean for a treated person}{$\mu_0(X_i)$}
\notearrow{blue}{Observed-control counterpart}{$\mu_0(X_{j(i)})$}
\notearrow{blue}{Correct the centering discrepancy}{Subtract $\hat\mu_0(X_i)-\hat\mu_0(X_{j(i)})$}
\notearrow{black}{More features: larger matching distances}{Fewer observations in each cell}
\begin{reviewnote}{qualification}{B1-014}{Matching bias and the estimand}{Chapter1-23.pdf, pp. 11, 12}
The relevant bias is $B_N=N_1^{-1}\sum_{i:T_i=1}[\mu_0(X_i)-\mu_0(X_{j(i)})]$, so the centered bias is $\sqrt{N_1}B_N$; the scan's single-$i$ expression abbreviates this average. Nearest-neighbor consistency does not itself imply negligible root-$N$ bias. Bias correction needs suitable smoothness, overlap and consistent regression approximation. Sample-share weights target ATE; treatment-variance weights generally target a different average.
\par\smallskip\textit{Reference:} \href{https://matheusfacure.github.io/python-causality-handbook/10-Matching.html}{Handbook, Chapter 10}; \href{https://www.nber.org/papers/t0283}{Abadie--Imbens, bias-corrected matching}.
\end{reviewnote}
\section{Chapter 11. Propensity scores}
\sourcepages{Chapter1-23.pdf, pp. 13--14}
Another way to address confounding bias is to summarize the confounders by a number between zero and one. The graph is $X\to p(X)\to T\to Y$, with $X\to Y$. The page writes $(Y_1,Y_0)\perp T\mid p(X)$ next to a double arrow and the balancing identity
\[E[T\mid p(X),X]=E[T\mid p(X)].\]


\textbf{My example.} Take two individuals, one treated and one control. If their probabilities of treatment are the same, they are comparable, even if their $X$ values differ.
\begin{align*}
E[T\mid p(X),X]&=E[T\mid X]=p(X), &&\text{since }p(X)\text{ is a function of }X,\\
E[T\mid p(X)]&=E[E[T\mid p(X),X]\mid p(X)]\\
&=E[p(X)\mid p(X)]=p(X).
\end{align*}

\paragraph{Inverse propensity weighting (IPW).}
The conditional difference can be rewritten as
\begin{align*}
E[Y\mid X,T=1]-E[Y\mid X,T=0]
={}&p(X)E\!\left[\frac{Y}{p(X)}\,\middle|\,X,T=1\right]\\
&-\{1-p(X)\}E\!\left[\frac{Y}{1-p(X)}\,\middle|\,X,T=0\right].
\end{align*}
\inkblue{?}\quad\textbf{My understanding.} We observe treatment for individual $i$, but the probability that $i$ would be treated is only $p(X_i)$, determined by $X$. Give every observed treated individual weight $1/p(X_i)$ to correct for this. If $p(X_i)$ is small, that observation contributes little to the unweighted outcome distribution; increase its representation with the reciprocal weight. For $T=1$ and $T=0$, the weights are respectively $1/p(X)$ and $1/(1-p(X))$. This reconstructs a population in which everyone is treated ($Y_1$) and another in which everyone is untreated ($Y_0$), and ATE is their difference.

A weighted variance is written
\[\hat\sigma_w^2=\frac{\sum_{i=1}^n w_i(Y_i-\hat\mu)^2}{\sum_{i=1}^n w_i}.\]
But the estimated propensity $\hat p(X)$ itself contains error. Adding nonconfounding covariates carelessly can make the propensity estimator highly variable. Treated people may have a higher estimated probability of treatment than untreated people. The bottom warning says that propensity weights above 20\% can increase variance (the threshold is recorded as written).

\paragraph{Propensity-score matching.}
Replace matching distance with $p(X)$. The author comments, ``this is not very good.'' Better prediction of treatment does not necessarily improve causal estimation and can increase variance. If treated and untreated propensity-score distributions lack overlap, matching encounters problems.

\begin{center}
\begin{tikzpicture}[>=Stealth, font=\small, line width=.65pt]

\node (x) at (0,1.5) {$X$};\node (p) at (1.4,.8) {$p(X)$};\node[draw,circle] (t) at (1.4,-.3) {$T$};\node[draw,circle] (y) at (0,-1.5) {$Y$};\draw[->] (x)--(p);\draw[->] (p)--(t);\draw[->] (t)--(y);\draw[->] (x)--(y);
\node (top) at (6,1) {$(Y_0,Y_1)\perp T\mid p(X)$};\node (bot) at (6,-1) {$E[T\mid p(X),X]=E[T\mid p(X)]$};\draw[<->,double] (top)--(bot);

\end{tikzpicture}
\par{\small Source p. 13: propensity-score sketch and the original double arrow (qualified in the review).}
\end{center}
\par\Needspace{9\baselineskip}\smallskip\textbf{Original pointing relations.}\par\smallskip
\notearrow{black}{Balancing identity: expand both sides}{$E[T\mid p(X),X]=E[T\mid p(X)]=p(X)$}
\notearrow{black}{Weight the observed treatment group}{Construct the all-treated $Y_1$ population}
\notearrow{black}{Weight the observed control group}{Construct the all-untreated $Y_0$ population}
\begin{reviewnote}{original}{B1-015}{Balancing is not equivalent to ignorability}{Chapter1-23.pdf, pp. 13}
The double arrow in the scan is not a logical equivalence. The identity $E[T\mid e(X),X]=E[T\mid e(X)]$ holds for the true binary-treatment propensity $e(X)=P(T=1\mid X)$. Conditional exchangeability $(Y_0,Y_1)\perp T\mid X$ implies $(Y_0,Y_1)\perp T\mid e(X)$. Overlap is additionally needed to identify both treatment-arm means on the target population; it is not necessary for this conditional-independence statement itself. Equal propensity scores balance covariate distributions; they do not make two particular people's potential outcomes identical.
\par\smallskip\textit{Reference:} \href{https://matheusfacure.github.io/python-causality-handbook/11-Propensity-Score.html}{Handbook, Chapter 11}.
\end{reviewnote}

\begin{reviewnote}{qualification}{B1-016}{Weighted dispersion is not the IPW estimator variance}{Chapter1-23.pdf, pp. 13, 14}
The displayed $\sum w_i(Y_i-\hat\mu)^2/\sum w_i$ measures weighted outcome dispersion, not the sampling variance of an IPW ATE estimator. With known propensities define $\psi_i=T_iY_i/e_i-(1-T_i)Y_i/(1-e_i)$; an i.i.d. variance estimate for its sample mean is $s_\psi^2/n$. Estimated nuisances require suitable influence-function or other valid inference. The scan's 20\% remark is a heuristic with unspecified normalization, not a universal threshold.
\par\smallskip\textit{Reference:} \href{https://matheusfacure.github.io/python-causality-handbook/11-Propensity-Score.html}{Handbook, Chapter 11}.
\end{reviewnote}
\section{Chapter 12. Doubly robust estimation}
\sourcepages{Chapter1-23.pdf, p. 15}
Combine linear outcome regression and propensity weighting to estimate the treatment-control contrast conditional on $X$:
\begin{align*}
\widehat{\mathrm{ATE}}={}&\frac1N\sum_i\left[\frac{T_i\{Y_i-\hat\mu_1(X_i)\}}{\hat p(X_i)}+\hat\mu_1(X_i)\right]\\
&-\frac1N\sum_i\left[\frac{(1-T_i)\{Y_i-\hat\mu_0(X_i)\}}{1-\hat p(X_i)}+\hat\mu_0(X_i)\right].
\end{align*}
Here $\hat p(X)$ estimates the propensity score (for example, using logistic regression), and $\hat\mu_1(X)$ and $\hat\mu_0(X)$ estimate $E[Y\mid X,T=1]$ and $E[Y\mid X,T=0]$. \inkblue{The first average estimates $E[Y_1]$ and the second $E[Y_0]$. Only one of the propensity or outcome models needs to be correct.}

\begin{inkblock}{blue}
\textbf{My algebraic explanation.} For $E[Y_1]$,
\[\widehat E[Y_1]=\frac1N\sum_i\left[\frac{T_i\{Y_i-\hat\mu_1(X_i)\}}{\hat p(X_i)}+\hat\mu_1(X_i)\right].\]
If $\hat p$ is wrong but $\hat\mu_1$ is correct, the first term is zero for $T_i=0$ and has mean zero for $T_i=1$, leaving the regression prediction. The scan's shorthand writes $E[\hat\mu_1(X_i)]=Y_i$.

Alternatively,
\[\widehat E[Y_1]=\frac1N\sum_i\left[\frac{T_iY_i}{\hat p(X_i)}-\frac{T_i-\hat p(X_i)}{\hat p(X_i)}\hat\mu_1(X_i)\right].\]
If the outcome model is wrong but the propensity model is correct, $T_i-\hat p(X_i)$ has mean zero conditional on $X_i$, so the correction vanishes in expectation. \end{inkblock}
The final reminder says the theory and the data must be right; try to use DRE even when beginning with a single model.

\par\Needspace{9\baselineskip}\smallskip\textbf{Original pointing relations.}\par\smallskip
\notearrow{black}{Combine the two approaches}{Outcome regression + propensity weighting $\longrightarrow$ doubly robust estimator}
\notearrow{blue}{First bracket estimates $E[Y_1]$, second $E[Y_0]$}{ATE is their difference}
\notearrow{blue}{When propensity is wrong but outcome model correct}{Mean residual cancellation in the first representation}
\notearrow{red}{The starred condition in the original: $E[T_i-\hat p(X_i)]=0$}{The alternative representation; condition on $X_i$ for cancellation}
\begin{reviewnote}{original}{B1-017}{Doubly robust cancellation is conditional}{Chapter1-23.pdf, pp. 15}
The source's $E[\hat\mu_1(X_i)]=Y_i$ is not an identity. If $\mu_1$ is correct, $E[T(Y-\mu_1(X))/q(X)\mid X]=0$. If the propensity $e$ is correct, $E[TY/e-(T-e)m(X)/e\mid X]=\mu_1(X)$ for any fixed $m$. Apply both arms and average. This requires consistency, conditional exchangeability, overlap and appropriate estimation control; ``one outcome model correct'' means both treatment-arm outcome regressions for an ATE.
\par\smallskip\textit{Reference:} \href{https://matheusfacure.github.io/python-causality-handbook/12-Doubly-Robust-Estimation.html}{Handbook, Chapter 12}.
\end{reviewnote}
\section{Chapter 13. Difference-in-differences}
\sourcepages{Chapter1-23.pdf, p. 17; placed before Chapter 14 on pp. 16 and 18}
Difference-in-differences evaluates the effect of a macro-level intervention. $D$ indicates treatment and $T$ time; $Y_D(T)$ is the potential outcome under treatment $D$ in period $T$.

\paragraph{Two comparisons and their problems.}
The target among the treated is $E[Y_1(1)-Y_0(1)\mid D=1]$, comparing treatment and no treatment in the same period; $Y_0(1)$ is counterfactual. Before treatment, $E[Y_1(0)\mid D=1]=E[Y_0(0)\mid D=1]$. A before/after comparison uses the pre-intervention outcome to stand in for the post-intervention untreated outcome. This fails when there is a time trend: the two untreated potential outcomes need not coincide.

A treated/control comparison in the post period uses
\[E[Y(1)\mid D=1]-E[Y(1)\mid D=0].\]
It uses the control group's untreated post-period outcome to estimate the treated group's counterfactual. \inkblue{The blue annotation warns that this too can fail.}

Combine the two:
\begin{align*}
\widehat{\mathrm{ATET}}={}&\{E[Y(1)\mid D=1]-E[Y(1)\mid D=0]\}\\
&-\{E[Y(0)\mid D=1]-E[Y(0)\mid D=0]\}.
\end{align*}
\inkblue{The connecting arrows give the equivalent reading: the experimental group's post-minus-pre change minus the control group's post-minus-pre change. The margin notes a linear model with an interaction, an unknown variance, and ``check parallel [trends].''}

\par\Needspace{9\baselineskip}\smallskip\textbf{Original pointing relations.}\par\smallskip
\notearrow{blue}{Counterfactual in the same post period}{$Y_0(1)\mid D=1$}
\notearrow{blue}{Time comparison may fail}{Untreated potential outcomes can trend over time}
\notearrow{blue}{Alternative reading of the curved connecting arrows}{$(\text{treated post}-\text{treated pre})-(\text{control post}-\text{control pre})$}
\begin{reviewnote}{qualification}{B1-018}{Parallel trends is an identifying assumption}{Chapter1-23.pdf, pp. 16, 17, 18}
For the two-period ATT, require no anticipation and $E[Y_0(1)-Y_0(0)\mid D=1]=E[Y_0(1)-Y_0(0)\mid D=0]$, plus consistency, no spillovers and a comparable composition. This is a statement about untreated mean trends; it neither identifies each person's counterfactual nor requires full independence for both $d=0,1$. Pre-period checks cannot prove future parallel trends. The within transformation removes additive time-invariant effects, but causal FE estimation still needs exogeneity (typically $E[e_{it}\mid X_{i1},\ldots,X_{iT},U_i]=0$).
\par\smallskip\textit{Reference:} \href{https://matheusfacure.github.io/python-causality-handbook/13-Difference-in-Differences.html}{Handbook, Chapter 13}; \href{https://matheusfacure.github.io/python-causality-handbook/14-Panel-Data-and-Fixed-Effects.html}{Handbook, Chapter 14}.
\end{reviewnote}
\section{Chapter 14. Panel data and fixed effects}
\sourcepages{Chapter1-23.pdf, pp. 16 and 18}
Panel data repeatedly observe the same person at several times, holding that person's identity fixed and avoiding confounding bias from fixed characteristics. The notes describe the fixed-effects model as a way to account for the person [a short ``IV'' abbreviation is written in the scan].

Recall DID. For the treated city, estimate its untreated outcome after intervention by its untreated outcome before intervention plus the untreated city's change over time:
\[Y_0(1)\mid D=1
=Y_0(0)\mid D=1+\{Y_0(1)\mid D=0-Y_0(0)\mid D=0\}.
\]

Then $\widehat{\mathrm{ATT}}=Y_1(1)\mid D=1-\widehat Y_0(1)\mid D=1$: the first is known, the second estimated. The example annotation mentions the effect of advertising in a DMA.

Parallel trends means treatment assignment is unrelated to growth in the potential outcome over time. The scan writes $(Y_d(t)-Y_d(t-1))\perp D$.
\[y_{it}=\beta X_{it}+\gamma U_i+e_{it}.\]
\inkblue{$y_{it}$ is the outcome for person $i$ at time $t$, $X_{it}$ a time-varying treatment/covariate, and $U_i$ a set of unobserved characteristics that do not vary with time.}

\paragraph{Limits of panel data (top margin on the continuation page).}
\begin{inkblock}{blue}
Panel data cannot solve reverse causality, where cause and effect reinforce one another; fixed effects cannot control this. They also cannot remove unobserved confounders that vary over time.
\end{inkblock}

\paragraph{Residualizing a linear panel model.}
For two groups of features, $X_1,X_2$,
\[\hat y=\hat\beta_1X_1+\hat\beta_2X_2.\]
The drawn feature matrix has observations as rows and features as columns.
\begin{enumerate}
\item Regress $y$ on the second feature group: $\hat y^*=\hat\gamma_1X_2$.
\item Regress $X_1$ on $X_2$: $\hat X_1=\hat\gamma_2X_2$.
\item Take residuals $\widetilde X_1=X_1-\hat X_1$ and $\widetilde y=y-\hat y^*$.
\item Regress the outcome residual on the feature residual, yielding the coefficient on $X_1$.
\end{enumerate}

\paragraph{Within transformation.}
Compute $\ddot Y_{it}=Y_{it}-\bar Y_i$ and $\ddot X_{it}=X_{it}-\bar X_i$, then regress $\ddot Y_{it}$ on $\ddot X_{it}$. Since $\bar U_i=U_i$, the fixed effect disappears:
\begin{align*}
Y_{it}&=\beta X_{it}+\gamma U_i+e_{it},\\
\bar Y_i&=\beta\bar X_i+\gamma\bar U_i+\bar e_i,\\
Y_{it}-\bar Y_i&=\beta(X_{it}-\bar X_i)+\gamma(U_i-\bar U_i)+(e_{it}-\bar e_i),\\
\ddot Y_{it}&=\beta\ddot X_{it}+\ddot e_{it}.
\end{align*}
\inkblue{This removes all unobserved characteristics that are constant over time.} For time effects, such as inflation, add one dummy for each period; time-varying factors may also matter. Group the data by person and sum the deviations: a zero total for the fixed characteristic expresses its lack of within-person change.

The sketch shows several parallel purchase-versus-cost lines. A fixed-effects model fits a line for each city and assumes a constant causal effect across conditions.

\begin{center}
\begin{tikzpicture}[>=Stealth, font=\small, line width=.65pt]

\draw[->] (0,0)--(4,0) node[right] {cost};\draw[->] (0,0)--(0,3) node[above] {purchase};\draw[blue] (.5,1.5)--(1.5,2.5);\draw[blue] (1.3,1.1)--(2.3,2.1);\draw[blue] (2,.6)--(3,1.6);\draw[blue] (2.6,.25)--(3.6,1.25);

\end{tikzpicture}
\par{\small Source p. 18: parallel within-city purchase--cost lines (blue).}
\end{center}
\par\Needspace{9\baselineskip}\smallskip\textbf{Original pointing relations.}\par\smallskip
\notearrow{black}{Repeated observations of the same person}{Hold person-specific characteristics fixed}
\notearrow{blue}{Person $i$ at time $t$: outcome}{$y_{it}$}
\notearrow{blue}{Treatment changes over time}{$X_{it}$}
\notearrow{blue}{Unobserved characteristics do not change over time}{$U_i$}
\notearrow{blue}{$\bar U_i=U_i$}{The within transformation removes the fixed effect}
\section{Chapter 15. Synthetic control}
\sourcepages{Chapter1-23.pdf, pp. 19--21}
Panel methods require disaggregated data. If only aggregate data are available, or no similar control unit can be found, combine several units into one similar unit.

There are $J+1$ units. Unit 1 receives treatment; units $2,\ldots,J+1$ are untreated and form the donor pool. $Y_{jt}^{I}$ and $Y_{jt}^{N}$ are the potential outcomes with and without intervention for unit $j$ at time $t$. Treatment starts at $T_0$. For $t>T_0$,
\[\tau_{1t}=Y_{1t}^{I}-Y_{1t}^{N}.\]
\inkblue{The treated outcome is factual and the untreated one counterfactual; the effect is a function of time.} For weights $W=(w_2,\ldots,w_{J+1})$,
\[\widehat Y_{1t}^{N}=\sum_{j=2}^{J+1}w_jY_{jt}.\]
\inkblue{Regression can be used to choose these weights: minimize the squared pre-intervention distance between the treated unit and the weighted donor average.}

\paragraph{Synthetic control as regression.}
Build a virtual California from donor states. The matrix has years as rows, states as columns, and cigarette sales as entries. Set California's series to $Y$, other states' series to $X$, and regress $Y\sim X$. The example weights include $-0.436$ and $-0.032$ (39 weights for 39 states). \inkblue{Their weighted sum is the virtual state, chosen to minimize its discrepancy from the treated unit.}

\begin{inkblock}{blue}
This easily overfits: the virtual state can almost coincide with California before intervention, then fluctuate strongly afterward and exceed California. The note mentions Ridge or Lasso regression. The handwritten size comparison reads ``$T$ large, $N$ large'' [its exact intended inequality is unclear].
\end{inkblock}

\paragraph{Avoiding extrapolation.}
Earlier overfitting is attributed to extrapolation: weights place observations outside the range of the data. Impose a convex-combination restriction. The synthetic control still has the form $\widehat Y_{1t}^{N}=\sum_{j=2}^{J+1}w_jY_{jt}$, but choose $W$ to minimize
\[\|X_1-X_0W\|_V=
\left\{\sum_{h=1}^{k}v_h\left(X_{h1}-\sum_{j=2}^{J+1}w_jX_{hj}\right)^2\right\}^{1/2},
\qquad w_j\ge0,\quad\sum_{j=2}^{J+1}w_j=1.
\]
The drawings contrast unconstrained extrapolation from donor points to a treated point outside their triangle with a constrained fit inside the donor convex hull.

$v_h$ reflects each variable's importance in matching treatment and synthetic control; different $v_h$ values yield different weights. Choices mentioned are standardizing each variable to mean zero and unit variance, or giving greater importance to variables that predict $Y$ better. The implementation here gives every variable equal importance: define a loss function, use SciPy quadratic optimization with weights summing to one, and constrain each weight to $[0,1]$. The resulting weights are sparse, small, and sum to one (the handwritten illustrative vector includes zeros and an approximately $0.3852$ entry).

\paragraph{Significance and placebo comparisons.}
Use a Fisher exact/permutation-style test. Refit synthetic control and estimate an effect for each of the 39 states. California gives the actual intervention estimate; the other states give placebo effects. If untreated states also appear to have large effects, that is concerning. Post-intervention variation is larger than pre-intervention variation; some units fit poorly even before intervention and should be excluded. Set a pre-period fit threshold using
\[\mathrm{MSE}=\frac1N\sum_t(Y_t-\widehat Y_t^{\mathrm{synth}})^2.\]
The stated test steps are to compute all state effects, then form
\[p=\frac1N\sum_j\mathbf1\{\hat\tau_{\mathrm{California}}>\hat\tau_j\},\]
with the reminder that California's effect is negative. \inkblue{A blue annotation says only one state's effect is higher than California's, gives $p=0.0286$, and concludes that California's estimated intervention effect is very extreme (effective).}

\begin{center}
\begin{tikzpicture}[>=Stealth, font=\small, line width=.65pt]

\foreach \x/\y in {0/2,2/2,1/.5}{\fill (\x,\y) circle(2pt);}\fill (3,0) circle(2pt) node[right,blue] {treated};\draw[->] (0,2)--(3,0);\draw[->] (2,2)--(3,0);\node at (1.3,-.55) {extrapolation};
\draw[dashed] (6,2)--(8,2)--(7,.5)--cycle;\foreach \x/\y in {6/2,8/2,7/.5}{\fill (\x,\y) circle(2pt);}\fill (9,0) circle(2pt) node[right,blue] {treated};\draw[->,blue] (9,0)--(7.5,1.2);\node[blue] at (7.2,-.55) {convex hull};

\end{tikzpicture}
\par{\small Source p. 20: extrapolation versus a convex-hull restriction (schematic source geometry).}
\end{center}
\par\Needspace{9\baselineskip}\smallskip\textbf{Original pointing relations.}\par\smallskip
\notearrow{blue}{Factual versus counterfactual}{$Y^I_{1t}$ versus $Y^N_{1t}$}
\notearrow{blue}{Effect is a function of time}{$\tau_{1t}$}
\notearrow{blue}{Find these weights by regression}{$\widehat Y^N_{1t}=\sum_jw_jY_{jt}$}
\notearrow{blue}{Nearly exact pre-fit can overfit}{The synthetic unit fluctuates after intervention}
\notearrow{black}{Repeat the procedure for every state}{California: actual effect; other states: placebo effects}
\notearrow{black}{Poor pre-intervention fits}{Set a pre-period MSE exclusion threshold}
\begin{reviewnote}{qualification}{B1-019}{Aggregate panels and synthetic-control identification}{Chapter1-23.pdf, pp. 19, 20}
A panel can consist of aggregate cities or states; individual-level data are not required. Synthetic control additionally needs untreated donors, no spillovers/anticipation and a credible stable relationship connecting pre- and post-period untreated outcomes. Convex weights restrict extrapolation outside the donor convex hull; they do not by themselves eliminate overfitting or guarantee causal identification. Read the unclear $T,N$ size annotation as a warning, not a proved sample-size rule.
\par\smallskip\textit{Reference:} \href{https://matheusfacure.github.io/python-causality-handbook/15-Synthetic-Control.html}{Handbook, Chapter 15}.
\end{reviewnote}

\begin{reviewnote}{qualification}{B1-020}{Placebo tail probability and its denominator}{Chapter1-23.pdf, pp. 21}
For a negative effect the relevant one-sided tail is $\hat\tau_j\le\hat\tau_{\rm CA}$, i.e. more negative, not ``higher.'' An exact uniform-assignment rank test uses $p=\#\{j:S_j\le S_{\rm obs}\}/J_{\rm eligible}$ including the observed assignment (or $(1+b)/(B+1)$ for sampled permutations). Exactness requires an assignment/exchangeability justification and a consistently applied statistic/filter. The scan does not specify enough to reconstruct $0.0286$; observational placebo ranks alone are not automatically Fisher-exact.
\par\smallskip\textit{Reference:} \href{https://matheusfacure.github.io/python-causality-handbook/15-Synthetic-Control.html}{Handbook, Chapter 15}.
\end{reviewnote}
\section{Chapter 16. Regression discontinuity design}
\sourcepages{Chapter1-23.pdf, pp. 22--24}
A discontinuity: what effect does it reveal? Let $R$ be the running variable, with treatment assigned discontinuously at cutoff $c$:
\[D_i=\mathbf1\{R_i>c\}.\]
\inkblue{Below $c$ treatment is zero; above it treatment is one.} Thus $Y_1$ is observed for $R>c$, and $Y_0$ for $R<c$. The figure draws smooth potential-outcome curves $Y_1,Y_0$, with the observed curve switching from $Y_0$ to $Y_1$ at $c$. The example plots age against all deaths / motor-vehicle deaths / suicide, looking for a discontinuity.

If potential-outcome means are smooth, their left and right limits agree, e.g.
\[\lim_{r\to c^-}E[Y_{1i}\mid R_i=r]=\lim_{r\to c^+}E[Y_{1i}\mid R_i=r].\]
The discontinuity in observed outcomes therefore identifies
\begin{align*}
&\lim_{r\to c^+}E[Y_i\mid R_i=r]-\lim_{r\to c^-}E[Y_i\mid R_i=r]\\
&\quad=\lim_{r\to c^+}E[Y_{1i}\mid R_i=r]-\lim_{r\to c^-}E[Y_{0i}\mid R_i=r]\\
&\quad=E[Y_{1i}-Y_{0i}\mid R_i=c].
\end{align*}
The source labels this a local average treatment effect. \inkblue{\textbf{My annotation.} This is a local random experiment, closest to an RCT.}

Estimate the treatment effect by a linear regression with different slopes on either side:
\[y_i=\beta_0+\beta_1r_i+\beta_2\mathbf1\{r_i>c\}+\beta_3\mathbf1\{r_i>c\}r_i.\]
\inkblue{The low-side line is the first part; the high-side line adds the indicator and interaction.} The next page identifies $\beta_0$ with the left limit, $\beta_0+\beta_2$ with the right limit, and $\beta_2$ with the cutoff effect.  Standard errors assess significance.

\paragraph{Kernel weighting.}
RDD relies heavily on linear regression's extrapolation. A fit can focus too much on distant observations and fit badly at the cutoff. Give observations nearer the threshold higher weight. The triangular kernel is
\[K(R,c,h)=\mathbf1\{|R-c|<h\}\left(1-\frac{|R-c|}{h}\right).\]
\inkblue{The indicator keeps only observations within $c\pm h$; the second factor decreases as distance from $c$ increases. Use WLS, weighted least squares.}

\paragraph{Fuzzy RDD.}
At the cutoff, treatment need not switch deterministically from zero to one: only its probability jumps. Use cutoff crossing as an IV; ignoring noncompliance biases the treatment estimate toward zero. The diagrams show a jump in treatment probability and a corresponding jump in outcome. The note assumes $T_{1i}>T_{0i}$ for all $i$ and no negative changes [the prose is abbreviated]. The Wald ratio is
\[
\frac{\lim_{r\to c^+}E[Y_i\mid R_i=r]-\lim_{r\to c^-}E[Y_i\mid R_i=r]}
{\lim_{r\to c^+}E[T_i\mid R_i=r]-\lim_{r\to c^-}E[T_i\mid R_i=r]}
=E[Y_{1i}-Y_{0i}\mid T_{1i}>T_{0i},R_i=c].
\]
The margin distinguishes simply estimating an RD discontinuity from IV estimation.

The continuation says a large $p$-value means nonsignificance and a smaller effect [author's shorthand, not an equivalence between significance and effect size]. The McCrary test checks bunching in the running-variable distribution: is there a spike near the cutoff?

\begin{center}
\begin{tikzpicture}[>=Stealth, font=\small, line width=.65pt]

\draw[->] (0,0)--(5,0) node[right] {$R$};\draw[->] (0,0)--(0,3) node[above] {$Y$};
\draw (0.2,.2) .. controls (1,2.5) and (2,2.2) .. (4.8,2.7) node[right] {$Y_1$};\draw (.2,.2) .. controls (1,1.2) and (3,.4) .. (4.8,2.7) node[right,yshift=-12pt] {$Y_0$};
\draw[red,thick] (.2,.2) .. controls (1,1.2) and (1.9,.8) .. (2.5,.95)--(2.5,2.25) .. controls (3,2.3) and (4,2.3) .. (4.8,2.7);\draw[red,dashed] (2.5,0)--(2.5,.95);\node[red,below] at (2.5,0) {$c$};
\begin{scope}[xshift=7cm]\draw[->] (0,0)--(4,0) node[right] {$R$};\draw[->] (0,0)--(0,2.7) node[above] {$E[T\mid R]$};\draw[dashed] (0,.3)--(4,1);\draw[dashed] (0,1.5)--(4,2.2);\draw[red,thick] (0,.3)--(2,.65)--(2,1.85)--(4,2.2);\draw[dotted] (2,0)--(2,.65);\node[below] at (2,0) {$c$};\end{scope}

\end{tikzpicture}
\par{\small Source pp. 22--23: black potential curves; red observed curve switches at the cutoff.}
\end{center}
\par\Needspace{9\baselineskip}\smallskip\textbf{Original pointing relations.}\par\smallskip
\notearrow{blue}{Local experiment: closest to an RCT}{$E[Y_1-Y_0\mid R=c]$}
\notearrow{blue}{Below-cutoff regression}{$\beta_0+\beta_1r_i$}
\notearrow{blue}{Above-cutoff regression adds}{$\beta_2+\beta_3r_i$}
\notearrow{blue}{Indicator: only use data inside $c\pm h$}{$\mathbf1\{|R-c|<h\}$}
\notearrow{blue}{Farther from the threshold gives less weight}{$1-|R-c|/h$}
\notearrow{black}{Treatment-probability jump as first stage}{Outcome jump divided by treatment jump: fuzzy-RD Wald ratio}
\begin{reviewnote}{original}{B1-021}{Center the running variable before reading the intercept}{Chapter1-23.pdf, pp. 22, 23}
For the formula as written, $m(c^-)=\beta_0+\beta_1c$ and $m(c^+)=\beta_0+\beta_1c+\beta_2+\beta_3c$. Thus the jump is $\beta_2+\beta_3c$. Set $r_i=R_i-c$ and use $1\{r_i>0\}$ to obtain the stated intercept interpretation and jump $\beta_2$.
\par\smallskip\textit{Reference:} \href{https://matheusfacure.github.io/python-causality-handbook/16-Regression-Discontinuity-Design.html}{Handbook, Chapter 16}.
\end{reviewnote}

\begin{reviewnote}{qualification}{B1-022}{RD locality, fuzzy compliance, and density tests}{Chapter1-23.pdf, pp. 22, 23, 24}
Sharp RD identifies the cutoff ATE under continuity of both potential-outcome means and support on both sides; literal local randomization is a stronger interpretation. Fuzzy RD additionally needs a nonzero treatment-probability jump, local exclusion/exogeneity and monotonicity. Its reduced form equals first-stage jump times complier LATE; this is a scaled estimand, not a universal claim that every naive estimate is biased toward zero. McCrary's test targets a discontinuity in density at the cutoff, not any nearby spike, and cannot establish all RD assumptions.
\par\smallskip\textit{Reference:} \href{https://matheusfacure.github.io/python-causality-handbook/16-Regression-Discontinuity-Design.html}{Handbook, Chapter 16}; \href{https://www.nber.org/papers/t0334}{McCrary, running-variable density test}.
\end{reviewnote}

\begin{reviewnote}{original}{B1-023}{Nonsignificance does not imply a small effect}{Chapter1-23.pdf, pp. 24}
The original top line's implication is invalid. A test statistic such as $\hat\tau/\mathrm{SE}(\hat\tau)$ depends on both magnitude and precision. A large estimate with a large SE can yield a large $p$-value. Report the estimate and interval; nonsignificance alone does not show an effect is small or absent.
\par\smallskip\textit{Reference:} \href{https://matheusfacure.github.io/python-causality-handbook/03-Stats-Review-The-Most-Dangerous-Equation.html}{Handbook, Chapter 3}.
\end{reviewnote}
\section{Part II. Applications}
\section{Chapter 17. Predictive Models 101}
\sourcepages{Chapter1-23.pdf, p. 25}
Machine learning estimates $E[Y\mid X]$, the expected outcome at a given $X$, producing $\widehat E[Y\mid X]$. An overly complex model overfits. The cross-validation diagram separates training data, test data, and unseen data. Distinguish predictions from policies. Estimators and maximum likelihood are noted.

\inkblue{$R^2$ evaluates a continuous-outcome model, describing how much outcome variation the model explains.} The written range is $(-\infty,1)$, with a reminder that it measures performance relative to the mean. \inkblue{As tree depth rises, training $R^2$ increases while test $R^2$ can decrease: overfitting.}

\texttt{pandas.qcut} divides model-predicted values into quantile bins, i.e. model segments. A return strategy can use a threshold to decide whether to trade or to lend, and how much money to commit.

\par\Needspace{9\baselineskip}\smallskip\textbf{Original pointing relations.}\par\smallskip
\notearrow{black}{ML learns from data}{$\widehat E[Y\mid X]$}
\notearrow{blue}{Depth increases; train fit improves, test fit worsens}{Overfitting}
\notearrow{black}{Decision threshold}{Lend or not; determine how much to commit}
\begin{reviewnote}{original}{B1-024}{Include perfect fit in the R-squared range}{Chapter1-23.pdf, pp. 25}
With a nonconstant outcome and $R^2=1-\mathrm{SSE}/\mathrm{SST}$, perfect prediction gives $R^2=1$, so the range is $(-\infty,1]$. In-sample OLS with an intercept gives $0\le R^2\le1$; out-of-sample $R^2$ can be negative. The zero-variance outcome requires a separate convention.
\par\smallskip\textit{Reference:} \href{https://matheusfacure.github.io/python-causality-handbook/17-Predictive-Models-101.html}{Handbook, Chapter 17}.
\end{reviewnote}
\section{Chapter 18. Heterogeneous treatment effects and personalization}
\sourcepages{Chapter1-23.pdf, p. 26}
Unlike ordinary ML, the aim is to intervene and predict the effect of intervention. Consider $E[Y\mid X,T]$, where $X$ consists of external characteristics we cannot control (e.g. sales volume, as written), and $T$ is the treatment (e.g. price). Optimize $\arg\max_T E[Y\mid X,T]$.

Move from ATE to CATE. For binary treatment, ATE is $E[Y_1-Y_0]$; for continuous treatment, the note uses $E[y'(t)]$, the derivative of the outcome with respect to treatment. The conditional versions are
\[E[Y_1-Y_0\mid X],\qquad E[y'(t)\mid X].\]
\begin{inkblock}{blue}
\textbf{My annotation.} Add the condition $X$: calculate the expectation separately for each $X$, then choose a personalized policy.
\end{inkblock}

Predicting elasticity means estimating $\partial Y_i/\partial T_i$, approximated by
\[\frac{Y(T_i)-Y(T_i+\varepsilon)}{T_i-(T_i+\varepsilon)}.\]
But both outcomes are unobservable for the same individual: the causal-inference problem. In a simple regression $y_i=\beta_0+\beta_1t_i+\beta_2X_i+e_i$, the derivative is only $\beta_1$, a common elasticity. Add an interaction:
\[y_i=\beta_0+\beta_1t_i+\beta_2X_i+\beta_3t_iX_i+e_i,
\qquad \frac{\partial\hat y_i}{\partial t_i}=\hat\beta_1+\hat\beta_3X_i.
\]
\inkblue{The interaction makes CATE depend on the covariates.} \inkred{$\beta_1,\beta_3$ determine the derivative relationship.} The note also writes a finite difference $\{y(t+\varepsilon)-y(t)\}/\varepsilon$, followed by the practical one-unit approximation $\hat y(t+1)-\hat y(t)$.

\par\Needspace{9\baselineskip}\smallskip\textbf{Original pointing relations.}\par\smallskip
\notearrow{blue}{Condition on $X$: average separately}{Personalized policy}
\notearrow{red}{Interaction term}{$\beta_3t_iX_i$}
\notearrow{blue}{Covariates make the effect vary}{$\hat\beta_1+\hat\beta_3X_i$}
\begin{reviewnote}{qualification}{B1-025}{Causal conditional means and elasticity units}{Chapter1-23.pdf, pp. 26}
The policy target is $\arg\max_t E[Y(t)\mid X=x]$; it equals the observational conditional-mean optimization only under consistency, exchangeability and overlap. $X$ must be pre-treatment (e.g. past sales, not sales caused by today's price). A raw derivative is sensitivity; elasticity is $(t/y)\,\partial y/\partial t=\partial\log y/\partial\log t$. A one-unit difference equals a derivative only for a linear response or as an appropriate approximation.
\par\smallskip\textit{Reference:} \href{https://matheusfacure.github.io/python-causality-handbook/18-Heterogeneous-Treatment-Effects-and-Personalization.html}{Handbook, Chapter 18}.
\end{reviewnote}
\section{Chapter 19. Evaluating causal models}
\sourcepages{Chapter1-23.pdf, pp. 27--28}
The purpose is prediction and ranking (grouping). Use nonrandom data to estimate causal models. Compare:
\begin{enumerate}
\item $\mathrm{sales}_i=\beta_0+\beta_1\mathrm{price}_i+\beta_2X_i+\beta_3X_i\mathrm{price}_i+e_i$, a linear model (the feature terms can be matrices).
\item $\mathrm{sales}_i=G(X_i,\mathrm{price}_i)+e_i$, using a regression tree. Check $R^2$ for overfitting.
\item A random-number prediction as a placebo comparison.
\end{enumerate}
An elasticity model ranks units by treatment sensitivity. Within a segment, estimate elasticity by a univariate regression $y_i=\beta_0+\beta_1t_i+e_i$:
\[\hat\beta_1=\frac{\sum_i(t_i-\bar t)(y_i-\bar y)}{\sum_i(t_i-\bar t)^2},\]
where $\bar t$ is mean treatment and $\bar y$ mean outcome.

\paragraph{Cumulative elasticity curve.}
Sort by predicted elasticity, from most sensitive downward. Estimate the slope for the first group, then keep adding groups until all observations are included, where the curve reaches ATE:
\[\widehat{y'}(t)_k=\hat\beta_k=
\frac{\sum_{i=1}^{k}(t_i-\bar t)(y_i-\bar y)}{\sum_{i=1}^{k}(t_i-\bar t)^2}.\]
The goals in the note are that $\widehat{y'}(t)_k>\widehat{y'}(t)_{k+a}$ and that the difference be large. The sketch descends toward the overall effect and labels a well-performing ranking as preferable.

\paragraph{Cumulative gain curve.}
\[\widehat F(k)=\hat\beta_k\frac{k}{N}
=\frac{\sum_{i=1}^{k}(t_i-\bar t)(y_i-\bar y)}{\sum_{i=1}^{k}(t_i-\bar t)^2}\frac{k}{N}.\]
The sketch compares a curved gain line with a straight reference. For a confidence interval, the standard error is
\[S_{\hat\beta_1}=\sqrt{\frac{\sum_i\hat\varepsilon_i^2}{(n-2)\sum_i(t_i-\bar t)^2}}.\]

\begin{center}
\begin{tikzpicture}[>=Stealth, font=\small, line width=.65pt]

\draw[->] (0,0)--(4,0);\draw[->] (0,0)--(0,2.7);\draw (.2,2.5) .. controls (.8,1) and (2,.6) .. (3.8,.5);\draw[dashed] (0,.5)--(3.8,.5);\node[right] at (3.8,.5) {overall};\node[below] at (2,0) {ranked sample share};
\begin{scope}[xshift=7cm]\draw[->] (0,0)--(4,0);\draw[->] (0,0)--(0,2.7);\draw (0,0)--(3.6,2.1);\draw (0,0) .. controls (.8,2.2) and (2,2.6) .. (3.6,2.1);\node[below] at (2,0) {$k/N$};\end{scope}

\end{tikzpicture}
\par{\small Source pp. 27--28: cumulative slope (left) and gain against a straight reference (right).}
\end{center}
\par\Needspace{9\baselineskip}\smallskip\textbf{Original pointing relations.}\par\smallskip
\notearrow{black}{Feature terms can be matrices}{$\beta_2X_i$ and $\beta_3X_i\mathrm{price}_i$}
\notearrow{black}{Keep adding ranked groups}{At all observations, the cumulative slope reaches the overall effect}
\begin{reviewnote}{qualification}{B1-026}{Evaluate the ranking on an identified effect}{Chapter1-23.pdf, pp. 27, 28}
Fit the ranking separately from evaluation; subgroup univariate slopes are causal on a randomized evaluation sample (or with valid adjustment), not merely because the model was trained on observational data. At each $k$, recompute $\bar t_k,\bar y_k$ within the first $k$ observations. For negative price effects, ``most sensitive'' means more negative or larger absolute magnitude; the written decreasing-positive-curve convention needs the corresponding sign change. The displayed SE assumes i.i.d. homoskedastic linear errors and is pointwise, not simultaneous over the whole curve.
\par\smallskip\textit{Reference:} \href{https://matheusfacure.github.io/python-causality-handbook/19-Evaluating-Causal-Models.html}{Handbook, Chapter 19}.
\end{reviewnote}
\section{Chapter 20. Plug-and-play estimators}
\sourcepages{Chapter1-23.pdf, pp. 29--31}
The author marks this part as ``translation for reference.'' Previously, comparable treatment and control groups gave ATE. The source writes $Y(0),Y(1)\perp X$ here [the conditioning/assignment symbol is inconsistent with the surrounding discussion]. Now ask: for whom does treatment work better; who is more sensitive?
\[\tau_i=Y_i(1)-Y_i(0),\qquad
Y_i^{\mathrm{obs}}=\begin{cases}Y_i(1),&t=1,\\Y_i(0),&t=0,\end{cases}
\]
\[\mathrm{ATE}=\tau=E[Y_i(1)-Y_i(0)]=E[\tau_i],\qquad
\mathrm{CATE}=\tau(x)=E[Y_i(1)-Y_i(0)\mid X=x]=E[\tau_i\mid X=x].\]
With $y_i=\beta_0+\beta_1t_i+\beta_2X_i+\beta_3t_iX_i+e_i$, the fitted effect is $\hat\tau(x)=\hat\beta_1+\hat\beta_3x$. \inkblue{The blue note says the effect does not involve $\beta_2$, but if $X$ is unrelated to the outcome, $\beta_1,\beta_3$ will be biased.}

Outcome MSE is $E[(Y_i-\hat Y_i)^2]$. The aim is instead to minimize treatment-effect error, written in the source as
\[E[(\tau(x)-\hat\tau(x))^2]=E[(Y_i(1)-Y_i(0)-\hat\tau(x))^2].\]


\paragraph{Transformed outcomes for binary treatment.}
Introduce $Y_i^*=2Y_iT_i-2Y_i(1-T_i)$. Under 50/50 random assignment, $T\perp(Y(0),Y(1))$, so $E[TY(t)]=E[T]E[Y(t)]$. Also, $Y_iT_i=Y_i(1)T_i$ and $Y_i(1-T_i)=Y_i(0)(1-T_i)$. Therefore
\begin{align*}
E[Y_i^*\mid X_i=x]
&=2E[Y_i(1)T_i\mid X_i=x]-2E[Y_i(0)(1-T_i)\mid X_i=x]\\
&=2E[Y_i(1)\mid X_i=x]\,0.5-2E[Y_i(0)\mid X_i=x]\,0.5\\
&=\tau(x).
\end{align*}
For nonrandom assignment, introduce propensity $e(X_i)$, with the relevant conditional independence assumption, and use
\[Y_i^*=Y_i\frac{T_i-e(X_i)}{e(X_i)(1-e(X_i))}.\]
When treatment is randomized at a known probability $p$, substitute $p$ for $e(X_i)$. The derivation, underlined in red in the scan, expands its conditional expectation as
\begin{align*}
E[Y_i^*\mid X_i=x]
={}&\frac{E[Y_i(1)\mid X_i=x]E[T_i\mid X_i=x]}{e(x)}\\
&-\frac{E[Y_i(0)\mid X_i=x]E[1-T_i\mid X_i=x]}{1-e(x)}=\tau(x).
\end{align*}
From the cumulative gain curve, the transformed target provides an unbiased but very noisy estimate of an individual's treatment effect. With large data, this is less problematic.

\paragraph{Continuous treatment.}
The target is $\tau(x)=E[\partial Y_i(t)/\partial t\mid X=x]=E[\tau_i\mid X=x]$. The sketch shows a downward demand-versus-price slope. A linear slope estimate is
\[\hat\tau(x)=\frac{\operatorname{Cov}(Y_i,T_i)}{\operatorname{Var}(T_i)}
=\frac{\sum_i(Y_i-\bar Y)(T_i-\bar T)}{\sum_i(T_i-\bar T)^2}.
\]
For a single observation under randomized treatment, define
\[Y_i^*=(Y_i-\bar Y)\frac{T_i-\bar T}{\sigma_T^2}.\]
Set $V_i=(T_i-\bar T)/\sigma_T^2$. The derivation uses $E[V_i\mid X_i=x]=0$, $E[T_iV_i\mid X_i=x]=1$, and
\[E[T_i(T_i-\bar T)\mid X_i=x]=E[(T_i-\bar T)^2\mid X_i=x],\]
as well as $E[T_ie_i\mid X_i=x]=E[T_i\mid X_i=x]E[e_i\mid X_i=x]$. With $Y_i=\alpha+\beta T_i+e_i$ at $X_i=x$,
\begin{align*}
E[Y_i^*\mid X_i=x]
&=E[(Y_i-\bar Y)V_i\mid X_i=x]\\
&=E[(\alpha+\beta T_i+e_i-\bar Y)V_i\mid X_i=x]=\beta=\tau(x).
\end{align*}
The proposed continuous-treatment extension is written
\[Y_i^*=(Y_i-\bar Y)\frac{T_i-M(T_i)}{\{T_i-M(T_i)\}^2},\]
with a note that $M$ estimates $E[T_i\mid X_i=x]$.  \inkred{The author's red note says the denominator creates numerical difficulty; for the intended use, only relative ranking is needed, not a numerical explanation.}

For nonlinear relationships, the author writes ``no solution yet,'' then suggests $D_i=1/P_i^\alpha$, so $\log D=-\alpha\log P_i$.

\par\Needspace{9\baselineskip}\smallskip\textbf{Original pointing relations.}\par\smallskip
\notearrow{black}{CATE}{Add a treatment--covariate interaction}
\notearrow{blue}{The effect omits $\beta_2$ algebraically}{$\hat\tau(x)=\hat\beta_1+\hat\beta_3x$}
\notearrow{red}{Minimize treatment-effect risk}{$E[(\tau(X)-\hat\tau(X))^2]$}
\notearrow{red}{Conditional expectation expansion (red underline)}{$E[Y^*\mid X=x]=\tau(x)$ under the stated assumptions}
\notearrow{black}{Continuous-treatment target}{Demand-versus-price slope}
\notearrow{red}{Numerical difficulty; only relative ranking wanted}{The proposed denominator $\{T_i-M(T_i)\}^2$}
\begin{reviewnote}{original}{B1-027}{Assignment independence uses treatment, not covariates}{Chapter1-23.pdf, pp. 29}
The scan explicitly writes $(Y(0),Y(1))\perp X$. That does not make treatment groups comparable. Use $(Y(0),Y(1))\perp T$ for random assignment, or $(Y(0),Y(1))\perp T\mid X$ with overlap for adjusted observational identification. The CATE can vary with $X$.
\par\smallskip\textit{Reference:} \href{https://matheusfacure.github.io/python-causality-handbook/20-Plug-and-Play-Estimators.html}{Handbook, Chapter 20}.
\end{reviewnote}

\begin{reviewnote}{original}{B1-028}{CATE risk and individual-effect risk differ}{Chapter1-23.pdf, pp. 29}
Let $\Delta=Y(1)-Y(0)$ and $\tau(X)=E[\Delta\mid X]$. For a prediction $f(X)$ independent of evaluation noise, $E[(\Delta-f(X))^2]=E[(\tau(X)-f(X))^2]+E[\operatorname{Var}(\Delta\mid X)]$. The cross term vanishes by conditional centering. Thus the risks share a minimizer but are not equal unless the conditional effect variance is zero.
\par\smallskip\textit{Reference:} \href{https://matheusfacure.github.io/python-causality-handbook/20-Plug-and-Play-Estimators.html}{Handbook, Chapter 20}.
\end{reviewnote}

\begin{reviewnote}{qualification}{B1-029}{What transformed outcomes are unbiased for}{Chapter1-23.pdf, pp. 30}
Under conditional exchangeability and $0<e(X)<1$, $E[Y^*\mid X]=\tau(X)$; this is the precise conditional-mean guarantee. Design-unbiasedness for a fixed unit's individual effect additionally conditions on its potential outcomes under randomized assignment. Large $N$ does not cure a lack of overlap or infinite/heavy-tailed inverse-weight variance. The factorization in the displayed proof must hold conditional on $X$.
\par\smallskip\textit{Reference:} \href{https://matheusfacure.github.io/python-causality-handbook/20-Plug-and-Play-Estimators.html}{Handbook, Chapter 20}.
\end{reviewnote}

\begin{reviewnote}{qualification}{B1-030}{Continuous-treatment proof needs conditional moments}{Chapter1-23.pdf, pp. 31}
For $E[Y\mid T,X=x]=a(x)+\tau(x)T$, define $m(x)=E[T\mid X=x]$, $v(x)=\operatorname{Var}(T\mid X=x)>0$ and $\ell(x)=E[Y\mid X=x]$. Then $E[(Y-\ell(X))(T-m(X))/v(X)\mid X=x]=\tau(x)$. Global treatment moments suffice when the randomization distribution is the same for every $x$. Using in-sample $\bar T,\bar Y$ makes the written finite-sample expectation proof only approximate; use population moments, cross-fitting and consistency as appropriate.
\par\smallskip\textit{Reference:} \href{https://matheusfacure.github.io/python-causality-handbook/20-Plug-and-Play-Estimators.html}{Handbook, Chapter 20}.
\end{reviewnote}

\begin{reviewnote}{original}{B1-031}{A squared individual residual is not treatment variance}{Chapter1-23.pdf, pp. 31}
Both issues are in the handwritten formula. The nuisance input should be $X_i$, so $r_i=T_i-\hat m(X_i)$. Dividing by $r_i^2$ makes the proposed target $(Y_i-\bar Y)/r_i$, which can explode at zero and is not a generally unbiased CATE target. The conditional-moment target above uses $v(X_i)$ instead. Alternatively, $\widetilde Y_i/r_i$ with weight $r_i^2$ represents R-loss; retain its undivided form when $r_i=0$.
\par\smallskip\textit{Reference:} \href{https://matheusfacure.github.io/python-causality-handbook/20-Plug-and-Play-Estimators.html}{Handbook, Chapter 20}; \href{https://arxiv.org/abs/1712.04912}{Nie--Wager, R-learner}.
\end{reviewnote}
\begin{reviewnote}{original}{B1-043}{An irrelevant covariate does not automatically bias OLS}{Chapter1-23.pdf, pp. 29}
The word ``unrelated'' is legible in the original blue margin; the old English had reversed it. In a correctly specified OLS model with $E[e\mid T,X]=0$ and full-rank design, an irrelevant covariate does not itself bias $\hat\beta_1$ or $\hat\beta_3$: conditional on the design, $E[\hat\beta]=\beta$. It may reduce precision. Bias requires another mechanism, such as omitted confounding, misspecification, measurement error or regularization. The algebraic point that $\beta_2$ drops out of the treatment derivative remains correct.
\par\smallskip\textit{Reference:} \href{https://matheusfacure.github.io/python-causality-handbook/05-The-Unreasonable-Effectiveness-of-Linear-Regression.html}{Handbook, Chapter 5}; \href{https://matheusfacure.github.io/python-causality-handbook/20-Plug-and-Play-Estimators.html}{Handbook, Chapter 20}.
\end{reviewnote}

\section{Chapter 21. Meta-learners}
\sourcepages{Chapter1-23.pdf, pp. 32--33}
These are ML methods for CATE. The note says to fit on nonrandom data and validate on random data, then adds that one must still remove confounding. The targets are $E[Y_i(1)-Y_i(0)\mid X]$ or the conditional derivative for continuous treatment.

\paragraph{S-learner (``Go-Horse Learner'').}
Use one ML model $M_S$ to estimate $\mu(X,T)=E[Y\mid T,X]$. Fit with $(X,T)$, then pass the same $X$ through the model once with $T=0$ and once with $T=1$. The predicted outcome difference gives CATE. A cumulative gain curve can reveal overfitting; tune \texttt{max\_depth}.

The weakness is a tendency to shrink treatment effects toward zero. An S-learner is a regularized ML model, so regularization restricts the estimated treatment effect. If other covariates have strong predictive effects, it may drop the treatment variable entirely.

\paragraph{T-learner.}
Splitting the data addresses complete omission of treatment. Fit
\[\mu_0(x)=E[Y\mid T=0,X=x],\qquad \mu_1(x)=E[Y\mid T=1,X=x],\]
then estimate $\hat\tau(x)=M_1(x)-M_0(x)$. The diagram sends treatment and control data to separate models, then applies both models to the same $X$ and subtracts predictions.

\begin{inkblock}{blue}
\textbf{My annotation.} Splitting also shrinks the data sets. For example, abundant $T=0$ data may support a good fit, but scarce $T=1$ data force a simpler model to avoid overfitting. Regularization then creates a tension with effect estimation.
\end{inkblock}

\paragraph{X-learner.}
The first stage resembles the T-learner; another stage fits the imputed effects and uses a propensity model.
\begin{enumerate}
\item Fit $\hat M_0(x)\approx E[Y\mid T=0,X=x]$ and $\hat M_1(x)\approx E[Y\mid T=1,X=x]$. \inkred{The red note assumes the relationships can differ for $T=0$ and $T=1$.}
\item Impute treatment effects by differencing observed outcomes and predicted missing outcomes:
\[\hat\tau_i^{(0)}=\hat M_1(X_i)-Y_i\quad(T_i=0),\qquad
\hat\tau_i^{(1)}=Y_i-\hat M_0(X_i)\quad(T_i=1).
\]
\item Fit effect models $\hat M_{\tau0}(x)\approx E[\hat\tau^{(0)}\mid X=x,T=0]$ and $\hat M_{\tau1}(x)\approx E[\hat\tau^{(1)}\mid X=x,T=1]$.
\item Combine them:
\[\hat\tau(x)=\hat M_{\tau0}(x)\hat e(x)+\hat M_{\tau1}(x)\{1-\hat e(x)\}.\]
\end{enumerate}
The flowchart crosses the two outcome models over to the opposite groups, constructs CATE pseudo-outcomes, fits two further models, and combines them using propensity $PS(X)$. \inkred{The red annotation says a larger $e(X)$ means more $T=1$ observations and a more accurate corresponding outcome model, warranting greater weight.} This refers to the model supplying the imputed counterfactual used by the opposite group's effect regression.

\begin{center}
\begin{tikzpicture}[>=Stealth, font=\small, line width=.65pt]

\node (s0) at (0,0) {$(X,0)$};\node (s1) at (2,0) {$(X,1)$};\node[draw] (sm) at (1,-1) {$M_S$};\node (sy0) at (0,-2) {$\hat Y(0)$};\node (sy1) at (2,-2) {$\hat Y(1)$};\node (sc) at (1,-3) {CATE};\draw[->] (s0)--(sm);\draw[->] (s1)--(sm);\draw[->] (sm)--(sy0);\draw[->] (sm)--(sy1);\draw[->] (sy0)--(sc);\draw[->] (sy1)--(sc);
\node (tx) at (5,0) {$X$};\node[draw] (t0) at (4,-1) {$M_0$};\node[draw] (t1) at (6,-1) {$M_1$};\node (ty0) at (4,-2) {$\hat Y(0)$};\node (ty1) at (6,-2) {$\hat Y(1)$};\node (tc) at (5,-3) {CATE};\draw[->] (tx)--(t0);\draw[->] (tx)--(t1);\draw[->] (t0)--(ty0);\draw[->] (t1)--(ty1);\draw[->] (ty0)--(tc);\draw[->] (ty1)--(tc);
\node[draw] (xm0) at (9,0) {$M_0$};\node[draw] (xm1) at (11,0) {$M_1$};\node (xd0) at (9,-1.2) {$M_1(X)-Y_0$};\node (xd1) at (11.8,-1.2) {$Y_1-M_0(X)$};\node[draw] (xt0) at (9,-2.2) {$M_{\tau0}$};\node[draw] (xt1) at (11.8,-2.2) {$M_{\tau1}$};\node (xc) at (10.4,-3.4) {CATE};\draw[->] (xm1)--(xd0);\draw[->] (xm0)--(xd1);\draw[->] (xd0)--(xt0);\draw[->] (xd1)--(xt1);\draw[->] (xt0)--(xc);\draw[->] (xt1)--(xc);\node (ps) at (7.4,-3.4) {$PS(X)$};\draw[->] (ps)--(xc);

\end{tikzpicture}
\par{\small Source pp. 32--33: S-, T- and X-learner pointing relations.}
\end{center}
\par\Needspace{9\baselineskip}\smallskip\textbf{Original pointing relations.}\par\smallskip
\notearrow{black}{Remove confounding first}{Fit on nonrandom data; evaluate on random data}
\notearrow{blue}{Splitting shrinks each group}{Scarce treated data require stronger regularization}
\notearrow{blue}{Tension: keep the fit simple yet estimate the effect}{Avoiding overfit versus regularization bias}
\notearrow{red}{Actual outcome minus imputed missing outcome}{The two X-learner pseudo-effects}
\notearrow{red}{Larger $e(X)$: more treated data, more accurate $M_1$}{Greater weight on $M_{\tau0}(X)$}
\begin{reviewnote}{qualification}{B1-032}{Meta-learner identification and X-learner weights}{Chapter1-23.pdf, pp. 32, 33}
S-, T- and X-learners identify CATE only with consistency, conditional exchangeability, overlap and adequate estimation. In the X-learner, large $e(x)$ supports an accurate $M_1$ used to impute missing treated outcomes for controls; this motivates greater weight on $M_{\tau0}$. Choosing $g(x)=e(x)$ is a useful combination rule, not a universally variance-optimal one. Regularization can shrink effects, but its bias need not always point toward zero.
\par\smallskip\textit{Reference:} \href{https://arxiv.org/abs/1706.03461}{Künzel et al., meta-learners}.
\end{reviewnote}
\section{Chapter 22. Debiased / orthogonal machine learning}
\sourcepages{Chapter1-23.pdf, pp. 34--36}
The notebook cites ``Double Machine Learning for Treatment and Causal Parameters'' (2016). ML is used for nuisance parameters. Existing data are biased because features are correlated: for sales, these may include price, temperature, cost, and weekday, with complex relationships.

\paragraph{Frisch--Waugh--Lovell.}
For $\hat Y=\hat\beta_1X_1+\hat\beta_2X_2$, where $X_1,X_2$ are feature matrices (one feature per column, one observation per row) and the source calls $\hat\beta_1,\hat\beta_2$ row vectors, the same coefficient on $X_1$ is obtained by:
\begin{enumerate}
\item Regressing $Y$ on $X_2$, giving $\hat Y^*=\hat\gamma_1X_2$.
\item Regressing $X_1$ on $X_2$, giving $\hat X_1=\hat\gamma_2X_2$.
\item Taking $\widetilde X_1=X_1-\hat X_1$ and $\widetilde Y=Y-\hat Y^*$.
\item Regressing $\widetilde Y$ on $\widetilde X_1$ to recover $\hat\beta_1$.
\end{enumerate}
Thus $(Y-\widehat{Y\sim X})\sim(T-\widehat{T\sim X})$ estimates $\tau$ in
\[Y_i-E[Y_i\mid X_i]=\tau\{T_i-E[T_i\mid X_i]\}+\varepsilon_i.\]

\paragraph{``FWL on steroids'': ML for ATE.}
\[Y_i-\hat M_Y(X_i)=\tau\{T_i-\hat M_T(X_i)\}+\varepsilon_i.\]
\inkred{$\hat M_Y$ estimates $E[Y\mid X]$ and acts as a denoising model, explaining variation in $Y$ due to $X$.} \inkred{$\hat M_T$ estimates $E[T\mid X]$ and acts as a debiasing model, removing treatment-assignment bias associated with $X$.}
\begin{enumerate}
\item Estimate outcome $Y$ from features $X$ with flexible ML regression $M_Y$.
\item Estimate treatment $T$ from $X$ with $M_T$.
\item Obtain residuals $\widetilde Y=Y-M_Y(X)$ and $\widetilde T=T-M_T(X)$.
\item Regress outcome residuals on treatment residuals: $\widetilde Y=\alpha+\tau\widetilde T$.
\end{enumerate}
The note calls this ``two predictions, one regression''; $\tau$ is the causal ATE parameter.

\begin{inkblock}{blue}
\textbf{My overfitting explanation.} If $M_Y$ overfits, it captures more than the $Y$--$X$ relationship: it also absorbs some $Y$--$T$ relationship. The final effect is biased toward zero because the outcome model has swallowed the causal relationship instead of leaving it for residual regression. If $M_T$ overfits, it explains too much variation in $T$, leaving little treatment-residual variance. The final estimate then has high variance. The margin links this to an overly accurate or extreme propensity score and violation of positivity.
\end{inkblock}

To address overfitting, split the data into $K$ folds. For each fold, train ML models on the other $K-1$ folds and obtain residuals on the held-out fold. Combine predictions/residuals for all folds and fit $\widetilde Y=\alpha+\tau\widetilde T$.

\paragraph{CATE and the R-learner.}
Allow the slope to depend on $X$:
\[Y_i-M_Y(X_i)=\tau(X_i)\{T_i-M_T(X_i)\}+\varepsilon_i.\]
One final regression is $\widetilde Y_i=\alpha+\beta_1\widetilde T_i+\beta_2X_i\widetilde T_i+\varepsilon_i$. The final model estimates the sales effect by $M(\mathrm{Price}=1,X_i)-M(\mathrm{Price}=0,X_i)$. \inkred{The bottom red label is ``R-learner.''}

\paragraph{Nonparametric double/debiased ML.}
For a continuous relationship, such as price and sales, write
\[Y_i=\hat M_Y(X_i)+\tau(X_i)\{T_i-\hat M_T(X_i)\}+\hat\varepsilon_i.\]
Rearrange the residual:
\[\hat\varepsilon_i=\{Y_i-\hat M_Y(X_i)\}-\tau(X_i)\{T_i-\hat M_T(X_i)\}.\]
Minimizing its square gives the causal R-loss and estimates CATE:
\begin{align*}
\hat L_n(\tau)&=\frac1n\sum_{i=1}^n\left[\{Y_i-\hat M_Y(X_i)\}-\tau(X_i)\{T_i-\hat M_T(X_i)\}\right]^2\\
&=\frac1n\sum_{i=1}^n(\widetilde Y_i-\tau(X_i)\widetilde T_i)^2\\
&=\frac1n\sum_{i=1}^n\widetilde T_i^2\left(\frac{\widetilde Y_i}{\widetilde T_i}-\tau(X_i)\right)^2.
\end{align*}
\begin{inkblock}{blue}
\textbf{My interpretation.} This is predicting $\widetilde Y_i/\widetilde T_i$ with any prediction model, using weights $\widetilde T_i^2$.
\end{inkblock}

``What is nonparametric about it?'' distinguishes nonparametric from nonlinear. The underlying equation $\widetilde Y_i=\tau(X_i)\widetilde T_i+e_i$ is still linear in residual treatment. Double ML yields a local linear approximation to a nonlinear CATE relationship, finding a derivative at a particular $X_i$.

\subsection{An unfinished nonlinear extension}
\sourcepages{Chapter1-23.pdf, p. 38 main discussion, then facing-page supplement p. 37}
\textbf{Author's status: unfinished.} The heading is ``Non-Scientific Double/Debiased ML.'' The author wants counterfactual prediction (``CATE?'') and an entire outcome curve from individual data points. Replace
\[\widetilde Y_i=\tau(X_i)\widetilde T_i+e_i\]
by
\[\widetilde Y_i=\tau(X_i,\widetilde T_i)+e_i.\]
The note calls this an unusual function to estimate with ML: first predict $\widetilde T_i$, then use it and $X$ to fit $\widetilde Y$. It mentions any ML model, particularly LGBM. A handwritten prediction label reads ``price'' although the adjacent function is $\hat\tau(X_i,\widetilde T_i)$; that label is inconsistent with the outcome being modeled.

For a test set, follow the same processing steps as in training. The margin mentions subsequent cross-prediction and averaging model predictions into one value [the preceding short implementation abbreviation is unclear in the scan].

The flowchart uses two cross-fitted treatment models and two outcome models, obtains $T$ residuals and $Y$ residuals, and feeds these together with $X$ into a final model for CATE. Drawbacks: (1) regularization biases the effect toward zero; (2) behavior depends on the ML method. Boosted trees have difficulty making smooth predictions or extrapolating to previously unseen prices, and can return strange predictions there.

\paragraph{Facing-page step-by-step supplement.}
Let $Y$ be sales, $T$ price, and $X$ covariates. Use five-fold LGBM cross-prediction to estimate $\hat Y\sim X$ and $\hat T\sim X$. Compute $\widetilde Y=Y-\hat Y$ and $\widetilde T=T-\hat T$. Fit an ML model (LGBM is suggested) to $\widetilde Y$ using $(\widetilde T,X)$; call it \texttt{finalM}.

For prediction, obtain $\hat T'$ from the five LGBM models given $X$. The handwritten instruction says to subtract a proposed new treatment $T_{\mathrm{new}}$ from $\hat T'$ to form $\widetilde T'$. This proposed treatment is a counterfactual value. Pass $(\widetilde T',X)$ through the final model to get $\widetilde Y'$ for each $T_{\mathrm{new}}$. The note adds that zero can be supplied as a reference residual. Add the predicted baseline $\hat Y$ back to $\widetilde Y'$ to obtain the final outcome curve. This supplement explains the unfinished construction; it is not presented as a validated estimator.

\begin{center}
\begin{tikzpicture}[>=Stealth, font=\small, line width=.65pt]

\node (data) at (0,0) {$(X,T,Y)$};\node[draw,align=center] (mt) at (3,1) {cross-fit $M_T$\\predict $T$};\node[draw,align=center] (my) at (3,-1) {cross-fit $M_Y$\\predict $Y$};\node (tr) at (6,1) {$\widetilde T=T-\hat T$};\node (yr) at (6,-1) {$\widetilde Y=Y-\hat Y$};\node[draw,align=center] (final) at (9,0) {final learner\\$(X,\widetilde T)\mapsto\widetilde Y$};\node (out) at (9,-2) {``CATE?'' / outcome curve};\draw[->] (data)--(mt);\draw[->] (data)--(my);\draw[->] (mt)--(tr);\draw[->] (my)--(yr);\draw[->] (tr)--(final);\draw[->] (yr)--(final);\draw[->] (final)--(out);

\end{tikzpicture}
\par{\small Source pp. 37--38: cross-prediction, residuals, and the unfinished final learner.}
\end{center}
\par\Needspace{9\baselineskip}\smallskip\textbf{Original pointing relations.}\par\smallskip
\notearrow{blue}{Predict outcome}{$\hat M_Y(X)$ in the p.~36 residual equation}
\notearrow{blue}{Predict treatment}{$\hat M_T(X)$ in the p.~36 residual equation}
\notearrow{blue}{R-loss: what the R-learner minimizes}{$\hat L_n(\tau)$}
\notearrow{black}{Residual-on-residual regression}{Recover the coefficient on $X_1$ by FWL}
\notearrow{red}{Denoising model}{$\hat M_Y(X)$}
\notearrow{red}{Debiasing model}{$\hat M_T(X)$}
\notearrow{blue}{Outcome model overfits}{It absorbs treatment signal; residual effect can shrink}
\notearrow{blue}{Treatment model overfits}{Residual treatment variation shrinks; effect variance grows}
\notearrow{blue}{Prediction target and weights}{$\widetilde Y_i/\widetilde T_i$ with weight $\widetilde T_i^2$}
\notearrow{black}{Unfinished: replace the linear residual relationship}{$\widetilde Y_i=\tau(X_i,\widetilde T_i)+e_i$}
\notearrow{black}{Hypothetical counterfactual price}{$T_{\rm new}$ in the facing-page prediction procedure}
\begin{reviewnote}{qualification}{B1-033}{Residual regression is not always the ATE}{Chapter1-23.pdf, pp. 34, 35}
The partially linear model $Y=\theta T+g(X)+\varepsilon$, $E[\varepsilon\mid T,X]=0$, supports the constant causal slope. With heterogeneous binary-treatment effects, pooled residual OLS instead targets $E[e(X)(1-e(X))\tau(X)]/E[e(X)(1-e(X))]$, generally not ATE. Cross-fitting reduces own-observation overfitting; it does not replace exchangeability, overlap, moments or nuisance-rate conditions. For the partialling-out score, a sufficient remainder condition is $\|\hat m-m\|_2(\|\hat\ell-\ell\|_2+\|\hat m-m\|_2)=o_p(n^{-1/2})$. See DML Sections 4.1 and 5.1 for the distinction.
\par\smallskip\textit{Reference:} \href{https://arxiv.org/html/1608.00060v7#S4.SS1}{Chernozhukov et al., DML Section 4.1}; \href{https://arxiv.org/html/1608.00060v7#S5.SS1}{Chernozhukov et al., DML Section 5.1}.
\end{reviewnote}

\begin{reviewnote}{original}{B1-034}{R-loss alone does not identify a local derivative}{Chapter1-23.pdf, pp. 36}
This claim is present in the original notes and also overstates the textbook discussion. Minimizing the conditional R-loss yields $\tau^*(x)=\operatorname{Cov}(T,Y\mid X=x)/\operatorname{Var}(T\mid X=x)$. It equals the treatment derivative for a conditionally linear response, but not generally for nonlinear $E[Y\mid T,X]$. For example, $Y=T^2$ and $T\sim U(0,1)$ give slope $1$, while the derivative is $2t$. Nonparametric flexibility in $x$ does not by itself localize in treatment $t$.
\par\smallskip\textit{Reference:} \href{https://arxiv.org/html/1712.04912v4#S2}{Nie--Wager, Section 2 (R-loss)}; \href{https://matheusfacure.github.io/python-causality-handbook/22-Debiased-Orthogonal-Machine-Learning.html}{Handbook, Chapter 22}.
\end{reviewnote}

\begin{reviewnote}{original}{B1-035}{Prediction residual has the subtraction reversed}{Chapter1-23.pdf, pp. 37}
The p.~37 subtraction direction is a source error: training uses $\widetilde T=T-\hat m(X)$, so prediction must use $\widetilde T_{\rm new}=T_{\rm new}-\hat m(X)$ too. Reversing this sign changes the fitted relationship. On p.~38 the output labeled ``price'' should be the residual outcome prediction $\widehat{\widetilde Y}$, because the final learner's target is residual sales.
\par\smallskip\textit{Reference:} \href{https://matheusfacure.github.io/python-causality-handbook/22-Debiased-Orthogonal-Machine-Learning.html}{Handbook, Chapter 22}.
\end{reviewnote}

\begin{reviewnote}{qualification}{B1-036}{Keep the unfinished nonlinear idea as a proposal}{Chapter1-23.pdf, pp. 37, 38}
The author's tentative status is appropriate. For fixed functions $m,\ell$, the transformation $r=t-m(x)$ is invertible at fixed $x$: a consistent fit $f(x,r)=E[Y-\ell(X)\mid X=x,T-m(X)=r]$ reconstructs $E[Y\mid x,t]=\ell(x)+f(x,t-m(x))$. Calling this a causal response requires exchangeability and support; its derivative or two-dose contrast gives an effect. Arbitrary residual S-learning does not inherit the orthogonality/rate guarantees of the R-loss estimator, and unseen treatment values require unsupported extrapolation.
\par\smallskip\textit{Reference:} \href{https://matheusfacure.github.io/python-causality-handbook/22-Debiased-Orthogonal-Machine-Learning.html}{Handbook, Chapter 22}.
\end{reviewnote}
\begin{reviewnote}{original}{B1-041}{Matrix products must follow the stated layout}{Chapter1-23.pdf, pp. 18, 34}
The scan puts observations in rows of $X_j\in\mathbb R^{n\times p_j}$ and describes the coefficients as row vectors. With that convention, the conformable equation is $\hat Y=X_1\hat\beta_1^{\mathsf T}+X_2\hat\beta_2^{\mathsf T}$, not $\hat\beta_1X_1+\hat\beta_2X_2$. Alternatively adopt column coefficient vectors and write $X_1\hat\beta_1+X_2\hat\beta_2$. The fitted projections are $\hat Y^*=P_{X_2}Y$ and $\hat X_1=P_{X_2}X_1$. This repairs matrix orientation; the FWL residualization argument is unchanged.
\par\smallskip\textit{Reference:} \href{https://matheusfacure.github.io/python-causality-handbook/05-The-Unreasonable-Effectiveness-of-Linear-Regression.html}{Handbook, Chapter 5}.
\end{reviewnote}

\section{Chapter 23. Effect heterogeneity and nonlinearity}
\sourcepages{Chapter1-23.pdf, pp. 39--40}
Treatment effect heterogeneity (TEH) means treatment effects differ among individuals. The example studies a latent outcome:
\begin{align*}
Y_{\mathrm{latent}}&\sim N(-4.5+0.001\,\mathrm{income}+\mathrm{nudge}
+0.01\,\mathrm{nudge}\cdot\mathrm{age},\ 1),\\
\mathrm{nudge}&\sim B(0.5),\qquad \mathrm{age}\sim G(10,4).
\end{align*}
The income distribution is written $G(20,2)$. \inkblue{Income's additive coefficient is the same for everyone; age is the variable used to split heterogeneous treatment effects.} Define $\mathrm{conversion}=\mathbf1\{Y_{\mathrm{latent}}>0\}$, with ``50\%'' written beside the threshold. The stated true ATE is 1.4 on the latent scale, versus 0.394 for conversion. Conversion is not linear.

\paragraph{What the sketches show.}
For the effect on the latent outcome, the cumulative effect curve ranked by age starts high and declines toward ATE; the estimated-income ranking stays near ATE. For the effect on conversion, both age and estimated income show treatment-effect heterogeneity. Income is strongly predictive of conversion even though it does not interact with treatment on the latent scale.

A third sketch shows conversion probability as an S-shaped function of estimated income. Its lower end is labeled a low conversion rate, its upper end a high conversion rate. At both ends the curve is flatter and treatment effects are small. The final written comparison says high conversion / high baseline conversion gives a low treatment effect, followed by ``low conversion, high baseline conversion, high treatment effect.''

\paragraph{A second example: price response.}
The sales-versus-price sketch has two downward-sloping nonlinear curves for types $a$ and $b$. Type $a$ is labeled wealthy. Across increasing price intervals the drawn declines are roughly 1, 2, 3, 4 for $a$ and 2, 4, 5 for $b$. The effect depends both on type ($a$ or $b$) and on price, so ranking by one fixed elasticity is problematic. The last suggestion is to transform price to make the relationship with $Y$ linear, with an example written $-(\mathrm{price}^{4})$; the exponent is not fully clear.

\begin{center}
\begin{tikzpicture}[>=Stealth, font=\small, line width=.65pt]

\draw[->] (0,0)--(4,0) node[right] {percentile};\draw[->] (0,0)--(0,2.8);\draw (.3,2.5) .. controls (.7,.5) and (2,.4) .. (3.8,.3) node[right] {age};\draw (.1,.25) .. controls (1,.55) and (2,.2) .. (3.8,.3);\node at (2,1.9) {latent};\node at (1.4,-.4) {income near ATE};
\begin{scope}[xshift=7cm]\draw[->] (0,0)--(4,0) node[right] {percentile};\draw[->] (0,0)--(0,2.8);\draw (.2,2.5) .. controls (1,1.9) and (2,1.8) .. (3.8,1.7) node[right] {age};\draw (.2,.2) .. controls (1.1,1.6) and (2.5,2.1) .. (3.8,1.7);\node at (2,.55) {income};\node at (2,2.7) {conversion};\end{scope}
\begin{scope}[yshift=-5cm]\draw[->] (0,0)--(4,0) node[right] {income};\draw[->] (0,0)--(0,3) node[above] {conversion};\draw (.1,.15) .. controls (2,.15) and (1.4,2.6) .. (3.8,2.6);\node[left] at (0,.2) {low};\node[left] at (0,2.6) {high};\end{scope}
\begin{scope}[xshift=7cm,yshift=-5cm]\draw[->] (0,0)--(4.3,0) node[right] {price};\draw[->] (0,0)--(0,3) node[above] {sales};\draw (.2,2.7) .. controls (2.2,2.7) and (2.8,2.1) .. (4,1.5) node[right] {$a$: wealthy};\draw (.2,2.1) .. controls (2,1.9) and (2.6,.6) .. (4,0) node[above right,yshift=3pt] {$b$};\draw[dashed] (1.2,0)--(1.2,2.65);\draw[dashed] (2.7,0)--(2.7,2.2);\node at (.7,2.9) {1};\node at (1.8,2.8) {2};\node at (2.5,2.55) {3};\node at (3.4,2.05) {4};\node at (.7,2.2) {2};\node at (1.8,1.7) {4};\node at (3,.6) {5};\end{scope}

\end{tikzpicture}
\par{\small Source pp. 39--40: latent versus conversion rankings, saturation, and type/price dependence.}
\end{center}
\par\Needspace{9\baselineskip}\smallskip\textbf{Original pointing relations.}\par\smallskip
\notearrow{blue}{Same additive effect for every individual}{$0.001\,\mathrm{income}$ in the latent equation}
\notearrow{blue}{Variable for treatment-effect segmentation}{$\mathrm{age}$ in $0.01\,\mathrm{nudge}\cdot\mathrm{age}$}
\notearrow{black}{Income can show treatment-effect heterogeneity}{Conversion-scale effect curve}
\notearrow{black}{Treatment effect depends on both type and price}{The $a,b$ sales curves; one fixed ranking is problematic}
\begin{reviewnote}{original}{B1-037}{The written simulation does not produce the quoted conversion effect}{Chapter1-23.pdf, pp. 39}
This is an original-note inconsistency traceable to the textbook's prose/code mismatch. With shape--scale Gamma notation, the written income has mean $40$, not $4000$. The textbook code uses $\mathrm{income}=100G(20,2)$ and $\mathbf1\{Y_{\rm latent}>0.1\}$; its reported $0.39477$ is a simulated treated-minus-control difference. The latent ATE is $1+0.01E[\mathrm{age}]=1.4$. Under the unscaled handwritten DGP, the baseline latent mean is about $-4.46$, so neither roughly 50\% conversion nor the quoted conversion effect follows.
\par\smallskip\textit{Reference:} \href{https://matheusfacure.github.io/python-causality-handbook/23-Challenges-with-Effect-Heterogeneity-and-Nonlinearity.html}{Handbook, Chapter 23}.
\end{reviewnote}

\begin{reviewnote}{original}{B1-038}{Low baseline conversion does not always mean a large effect}{Chapter1-23.pdf, pp. 39}
The repeated ``high baseline conversion'' and the blanket low-conversion/high-effect conclusion are in the scan. For latent mean $a$, treatment shift $\delta>0$, unit-variance normal noise and threshold $c$, the conversion effect is $\Phi(a+\delta-c)-\Phi(a-c)$. It becomes small as $a\to-\infty$ or $a\to+\infty$ and is largest in an intermediate region. The author's preceding ``both ends are flat'' intuition is the consistent one.
\par\smallskip\textit{Reference:} \href{https://matheusfacure.github.io/python-causality-handbook/23-Challenges-with-Effect-Heterogeneity-and-Nonlinearity.html}{Handbook, Chapter 23}.
\end{reviewnote}

\begin{reviewnote}{qualification}{B1-039}{A power transform is a model proposal}{Chapter1-23.pdf, pp. 40}
The price transformation is a tentative modeling idea; the handwritten exponent is not fully resolved. A transformation $q=h(p)$ linearizes a curve only if its functional form supports that choice. If $E[Y\mid p,x]=a(x)+b(x)h(p)$, the price effect remains $b(x)h'(p)$, so a constant slope in transformed price need not imply a constant price effect. Preserve the original $a/b$ and price-dependent ranking insight.
\par\smallskip\textit{Reference:} \href{https://matheusfacure.github.io/python-causality-handbook/23-Challenges-with-Effect-Heterogeneity-and-Nonlinearity.html}{Handbook, Chapter 23}.
\end{reviewnote}

}
